\documentclass{article}
\usepackage{anyfontsize}
\usepackage[UTF8]{ctex} % 如果需要支持中文
\usepackage{fontspec} 

\usepackage[a4paper, left=0.1in, right=0.1in, top=0.05in, bottom=0.05in]{geometry} % 设置四边页边距为0.1英寸{geometry} % 设置页边距为0.5英寸
\usepackage{enumitem} % 用于自定义列表
\usepackage{amsmath}  % 数学公式支持

\setlength{\parindent}{0pt} % 不缩进
\usepackage{etoolbox} % 用于列表操作
%调整类代码
\usepackage{comment}
\usepackage{tocloft} % 用于定制目录样式
\usepackage{hyperref} %不需要目录盒子注释这
\usepackage{ifthen}
\usepackage{tikz}
\usepackage{titletoc}
\usepackage{graphicx}
\usepackage{amssymb}  % 提供数学符号，包括\therefore
%目录设定
%快捷键 CMD + / 注释
%  \renewcommand{\cftdot}{} % 隐藏目录中的页码
%  \cftpagenumbersoff{section}
%  \cftpagenumbersoff{subsection}
%  \makeatletter
%  \renewcommand{\@pnumwidth}{0em}  % 设置页码宽度为0
%  \renewcommand{\@tocrmarg}{0em}   % 设置右边距为0
%  \makeatother
 
%\setCJKmainfont[
 %   Path = C:/USERS/11252/APPDATA/LOCAL/MICROSOFT/WINDOWS/FONTS/,
 %   UprightFont = {SourceHanSerifCN-Light-5.otf},
 %   BoldFont = {SourceHanSerifCN-Medium-6.otf},
%    Scale=1
%]{SourceHanSerifCN}

%\setmainfont{SourceHanSerif}  % 设置英文字体

% 处理冲突的命令
\let\oldbackepsilon\backepsilon
\let\backepsilon\relax
\let\oldeth\eth
\let\eth\relax
\let\olddigamma\digamma
\let\digamma\relax

\usepackage{float}
\usepackage[a4paper, left=0.1in, right=0.1in, top=0.05in, bottom=0.05in]{geometry} % 设置四边页边距为0.1英寸{geometry} % 设置页边距为0.5英寸
\usepackage{enumitem} % 用于自定义列表
\usepackage{amsmath}  % 数学公式支持

\setlength{\parindent}{0pt} % 不缩进
\usepackage{etoolbox} % 用于列表操作
%调整类代码
\usepackage{comment}
\usepackage{tocloft} % 用于定制目录 样式
\usepackage{hyperref} %不需要目录盒子注释这
\usepackage{ifthen}
\usepackage{tikz}
\usepackage{titletoc}
\usepackage{graphicx}
\usepackage{amssymb}  % 提供数学符号，包括\therefore
%目录设定
%快捷键 CMD + / 注释
%  \renewcommand{\cftdot}{} % 隐藏目录中的页码
%  \cftpagenumbersoff{section}
%  \cftpagenumbersoff{subsection}
%  \makeatletter
%  \renewcommand{\@pnumwidth}{0em}  % 设置页码宽度为0
%  \renewcommand{\@tocrmarg}{0em}   % 设置右边距为0
%  \makeatother
 


\usepackage{environ}

\newif\ifshowcontent  % 定义一个条件判断变量
\showcontenttrue    % 设置为 false，隐藏所有 question 环境的内容

\newcounter{question} % 题目计数器
\setcounter{question}{0}
\NewEnviron{question}{%
    \ifshowcontent
        \par\stepcounter{question}\textbf{\thequestion.} \BODY\par % 如果显示内容，则输出
    \fi
}

\NewEnviron{choices}{%
    \ifshowcontent
        \begin{enumerate}[label=\Alph*., leftmargin=*]
            \BODY
        \end{enumerate}\par
    \fi
}

\NewEnviron{correctanswer}{%
    \ifshowcontent
        \textbf{答案:}\BODY\par
    \fi
}



\newenvironment{explanation}
{\textbf{解析:}\par}
{\par}



% 新增考察方面命令
\newcommand{\checklist}[1]{
    \textbf{考察:} #1 \par % 在题目下方显示考察方面
    \addcontentsline{toc}{subsection}{题号\thequestion 考察: #1} % 将考察内容添加到目录
    \vspace{2em} % 添加1em的垂直空间
}

\graphicspath{{images/}} %统一


\begin{document}
 %\fontsize{22}{31}\selectfont
 \tableofcontents % 添加目录
\newpage
%\pagestyle{empty}




\section{考点1:市场风险度量}
\setcounter{question}{0}


\begin{question}
    Which of the following statements concerning expected shortfall estimates and coherent risk measures are correct?
\end{question}

\begin{choices}
    \item  Expected shortfall and coherent risk measures could estimate quantiles for the entire loss distribution both.
    \item  Expected shortfall and coherent risk measures could estimate quantiles for the tail region both.
    \item  Expected shortfall could estimate quantiles for the tail region and coherent risk measures could estimate quantiles for the non-tail region only.
    \item  Expected shortfall could estimate quantiles for the entire distribution and coherent risk measures could estimate quantiles for the tail region only.
\end{choices}
\begin{correctanswer}
    B
\end{correctanswer}
\begin{explanation}
    \textbf{A选项错误。}
    \begin{itemize}
        \item 选项A认为两者都估计整个分布的分位数。但ES仅关注尾部，而一致性风险度量可以覆盖整个分布。
    \end{itemize}
    \textbf{B选项正确。}
    \begin{itemize}
        \item ES对\textbf{尾部区域（tail region）}中n-1相同权重的分位数进行估算。
        
        一致性风险度量（coherent risk measure）估值包括尾部区域在内的\textbf{整个分布}的分位数。
    \end{itemize}
    \textbf{C选项错误。}
    \begin{itemize}
        \item 选项C认为ES估计尾部，而一致性风险度量仅估计非尾部，但一致性风险度量可以涵盖整个分布。
    \end{itemize}
    \textbf{D选项错误。}
    \begin{itemize}
        \item 选项D认为ES估计整个分布，而一致性风险度量仅尾部，与事实相反。
    \end{itemize}
\end{explanation}
\checklist{预期损失（Expected shortfall）和一致性风险度量（Coherent risk measure）的对比}


\begin{question}
    The VaR at a 95\% confidence level is estimated to be 1.732 from a historical simulation of 1,000 observations. Which of the following statements is most likely correct?
    
\end{question}

\begin{choices}
    \item The parametric assumption of normal returns is correct.
    \item The parametric assumption of lognormal returns is correct.
    \item The historical distribution has fatter tails than a normal distribution.
    \item The historical distribution has thinner tails than a normal distribution.
\end{choices}



\begin{correctanswer}
    C
\end{correctanswer}

\begin{explanation}
    \textbf{A选项错误。}
    \begin{itemize}
        \item 若正态假设正确，VaR 应接近 1.645，而非 1.732。

        当前值远高于 1.645，说明偏离正态。
    \end{itemize}
    \textbf{B选项错误。}
    \begin{itemize}
        \item 对数正态通常用于资产价格而不是收益率，而VaR一般是以收益率或损失角度计算的。

        此外没有证据支持 lognormal 分布。
    \end{itemize}
    \textbf{C选项正确。}
    \begin{itemize}
        \item 正态分布下，95\% VaR 对应的 Z 值是 1.645。

        但历史模拟法得到的VaR为 1.732，高于1.645。
        
        这说明在历史数据中，尾部的极端损失比正态分布更严重，即“尾部更厚”或“更重尾”（fatter tails）。
        
        这意味着如果只用正态分布建模风险将被低估。
    \end{itemize}
    \textbf{D选项错误。}
    \begin{itemize}
        \item 完全相反。如果是更轻尾，那 VaR 应该比 1.645 更小，而不是更大。
    \end{itemize}
\end{explanation}
\checklist{理解VaR的意义}


\begin{question}
    What is the difference between using a 95\% value at risk and a 95\% expected shortfall for a bond portfolio with \$825 million in assets and a probability of default of 3\%?
    
\end{question}

\begin{choices} 
    \item Both measures will show the same result.
    \item The VaR shows a loss of \$495 million while the ES shows no loss.
    \item The VaR shows no loss while the ES shows a \$495 million loss.
    \item The VaR shows no loss while the ES shows a \$395 million loss.
\end{choices}

\begin{correctanswer}
    C
\end{correctanswer}

\begin{explanation}
    
    \textbf{C选项正确。}
    \begin{itemize}
        \item The VaR measure would show a \textdollar 0 loss because the probability of default is less than 5\%. Having a 3\% probability means that three out of five times, in the tail, the portfolio will experience a total loss. The potential loss is \textdollar 495 million(= 3/5 x \textdollar 825 million).
        \item VaR显示\textdollar 0亏损，因为违约的概率小于5\%。 3\%的概率意味着，在尾部时，投资组合有五分之三的概率会全部损失。潜在损失为 3/5 x 8.25亿美元=4.95亿美元。
    \end{itemize}
  
\end{explanation}
\checklist{VaR和ES的差异}


\begin{question}
    A portfolio has a current value of \$40 million and its geometric returns are normally distributed with annual mean of 12.0\% and annual volatility (standard deviation) of 40\%. Because the geometric returns are normal, the future value of the portfolio has a lognormal distribution. If the returns are i.i.d., what is the 95\% confident level, 10-day lognormal value at risk?
\end{question}

\begin{choices}
    \item \$2.91 million
    \item \$3.86 million
    \item \$4.77 million
    \item \$5.07 million
\end{choices}

\begin{correctanswer}
    C
\end{correctanswer}

\begin{explanation}
    \textbf{C选项正确。}
    \begin{itemize}
        \item 10-day mean $= 12\% \times \dfrac{10}{250} = 0.48\%$
        \item $\text{10-day volatility} = 40\% \times \sqrt{\dfrac{10}{250}} = 8\% $
        \item $\text{Lognormal VaR} = 40 \times \left(1 - e^{0.48\% - 8\% \times 1.645}\right) = 4.76$
            
    \end{itemize}

\end{explanation}
\checklist{计算多期VaR}


\begin{question}
    A sovereign wealth fund has USD 50 billion in assets. The risk manager computes the daily VaR at
    various confidence levels as follow:
    \begin{table}[H]
        \centering
        \begin{tabular}{|c|c|}
            \hline
            Confidence Level & VaR(USD) \\
            \hline
            95.5\% & 720,000,000 \\
            96.0\% & 750,000,000 \\
            96.5\% & 790,000,000 \\
            97.0\% & 830,000,000 \\
            97.5\% & 870,000,000 \\
            98.0\% & 920,000,000 \\
            98.5\% & 990,000,000 \\
            99.0\% & 1,050,000,000 \\
            99.5\% & 1,150,000,000 \\
            \hline
        \end{tabular}
    \end{table}

    What is the closest estimate of the daily expected shortfall at the 97.5\% confidence level?
\end{question}

\begin{choices}
    \item USD 820 million
    \item USD 870 million
    \item USD 920 million
    \item USD 1,028 million
\end{choices}

\begin{correctanswer}
    D
\end{correctanswer}

\begin{explanation}
   
    \textbf{选项D正确。}
    
    97.5\%置信水平下的预期损失通过计算97.5\%以上置信水平VaR的平均值得到：
    (920,000,000 + 990,000,000 + 1,050,000,000 + 1,150,000,000)/4 = 1,027,500,000
    
    因此，预期损失约为1,028百万美元。
    
    \textbf{选项A错误。} 该数值过低，不代表尾部VaR的平均值。
    
    \textbf{选项B错误。} 这是97.5\%置信水平下的VaR值，而非预期损失。
    
    \textbf{选项C错误。} 这是98.0\%置信水平下的VaR值，而非尾部VaR的平均值。
\end{explanation}
\checklist{预期损失计算：取尾部VaR平均值}


\begin{question}
    Which of the following statements comparing VaR with expected shortfall is true?
\end{question}

\begin{choices}
    \item  Expected shortfall is sub-additive while VaR is not.
    \item  Both VaR and expected shortfall measure the amount of capital an investor can expect to lose over a given time period and are, therefore, interchangeable as risk measures.
    \item  Both VaR and expected shortfall depend on the assumption of a normal distribution of returns.
    \item  VaR can vary according to the confidence level selected, but expected shortfall will not.
\end{choices}

\begin{correctanswer}
    A
\end{correctanswer}

\begin{explanation}
    \textbf{选项A正确。}
    
    预期损失具有次可加性，而VaR不具有次可加性。这意味着当两个投资组合合并时，预期损失会小于或等于各自预期损失之和，而VaR可能违反这一性质。
    
    \textbf{选项B错误。} VaR和预期损失虽然都衡量风险，但衡量方式不同，不能互换使用。VaR衡量给定置信水平下的最大可能损失，而预期损失衡量超过VaR的预期损失。
    
    \textbf{选项C错误。} VaR和预期损失都可以基于正态分布假设，但也可以使用历史模拟法等其他方法，不一定依赖于正态分布假设。
    
    \textbf{选项D错误。} 改变置信水平时，VaR和预期损失都会发生变化，因为预期损失是基于VaR计算的。
\end{explanation}
\checklist{VaR与ES的次可加性差异}

\begin{question}
    The annual mean and volatility of a portfolio are 15\% and 40\%, respectively. The current value of the portfolio is USD 200,000. How does the 1-year 95\% VaR that is calculated using a normal distribution assumption (normal VaR) compare with the 1-year 95\% VaR that is calculated using the lognormal distribution assumption (lognormal VaR)?
\end{question}

\begin{choices}
    \item Lognormal VaR is greater than normal VaR by USD 22,400
    \item Lognormal VaR is greater than normal VaR by USD 28,600
    \item Lognormal VaR is less than normal VaR by USD 22,400
    \item Lognormal VaR is less than normal VaR by USD 28,600
\end{choices}

\begin{correctanswer}
    C
\end{correctanswer}

\begin{explanation}
    \textbf{选项C正确。}
    
    步骤1：计算正态分布VaR
    正态VaR = |0.15 - 1.645 × 0.40| = 0.508
    
    步骤2：计算对数正态分布VaR
    对数正态VaR = 1 - exp[0.15 - 1.645 × 0.40] = 0.398
    
    步骤3：计算差异
    差异 = 200,000 × (0.508 - 0.398) = 22,400
    
    因此，对数正态VaR比正态VaR小22,400美元。
    
    \textbf{选项A错误。} 对数正态VaR小于正态VaR，不是大于。
    
    \textbf{选项B错误。} 对数正态VaR小于正态VaR，且差异不是28,600美元。
    
    \textbf{选项D错误。} 差异是22,400美元，不是28,600美元。
\end{explanation}
\checklist{正态与对数正态分布VaR计算差异}

\begin{question}
    Given a VaR equal to 4.20, a threshold of 2.0\%, a shape parameter equal to 0.30, and a scale parameter equal to 0.45, what is the expected shortfall?
\end{question}

\begin{choices}
    \item 5.657
    \item 5.914
    \item 6.234
    \item 5.342
\end{choices}

\begin{correctanswer}
    A
\end{correctanswer}

\begin{explanation}
    \textbf{选项A正确。}
    
    使用极值理论(EVT)的预期损失公式：
    \[ \mathrm{ES} = \frac{\mathrm{VaR}}{1-\xi} + \frac{\beta - \xi u}{1-\xi} = \frac{4.20}{1-0.30} + \frac{0.45 - 0.30 \times 2.0}{1-0.30} = 5.657\]
    
    其中：VaR = 4.20，形状参数ξ = 0.30，尺度参数β = 0.45，阈值u = 2.0%
    
    \textbf{选项B错误。} 计算结果不正确，应为5.657。
    
    \textbf{选项C错误。} 计算结果不正确，应为5.657。
    
    \textbf{选项D错误。} 计算结果不正确，应为5.657。
\end{explanation}
\checklist{极值理论预期损失计算公式应用}

\begin{question}
    Which of the following combinations would most likely increase the standard errors of coherent risk measures?
\end{question}

\begin{choices}
    \item Larger sample size and higher tail probability
    \item Larger sample size and lower tail probability
    \item Smaller sample size and higher tail probability
    \item Smaller sample size and lower tail probability
\end{choices}

\begin{correctanswer}
    C
\end{correctanswer}

\begin{explanation}
    \textbf{选项C正确。}
    
    风险度量的标准误差公式为：
    \[
    se(q) = \frac{\sqrt{\dfrac{p(1-p)}{n}}}{f(q)}
    \]
    
    其中：n为样本量，p为尾部概率，f(q)为概率密度函数。
    
    较小的样本量(n)会直接增加标准误差。此外，当尾部概率(p)增加时，p(1-p)项会增加直到p=0.5，这也会增加标准误差。
    
    \textbf{选项A错误。} 较大的样本量会减少标准误差，而不是增加。
    
    \textbf{选项B错误。} 较大的样本量会减少标准误差，较低的尾部概率也会减少标准误差。
    
    \textbf{选项D错误。} 较小的样本量会增加标准误差，但较低的尾部概率会减少标准误差，综合效果不如选项C。
\end{explanation}
\checklist{样本量与尾部概率对风险度量误差的影响}


\begin{question}
    A credit union has accumulated historical loan returns data. The credit union believes the returns follow a normal distribution with a mean of 10 and a standard deviation of 3.0. The following table shows tail VaRs at different confidence levels. Assume the initial analysis uses four tail slices. Calculate the expected shortfall at the 95\% confidence level and identify the effect on ES when the number of tail slices increases.
    \begin{table}[H]
        \centering
        \begin{tabular}{|c|c|}
            \hline
            Confidence level & Tail VaR \\
            \hline
            95\% & 3.20 \\
            96\% & 3.50 \\
            97\% & 3.90 \\
            98\% & 4.40 \\
            99\% & 5.00 \\
            \hline
        \end{tabular}
    \end{table}
\end{question}

\begin{choices}
    \item (Expected Shortfall) 3.95; (Increasing Slices) ES increases
    \item (Expected Shortfall) 3.95; (Increasing Slices) ES decreases
    \item (Expected Shortfall) 4.20; (Increasing Slices) ES increases
    \item (Expected Shortfall) 4.20; (Increasing Slices) ES decreases
\end{choices}

\begin{correctanswer}
    C
\end{correctanswer}

\begin{explanation}
    \textbf{选项C正确。}
    
    预期损失通过计算95\%置信水平以上尾部VaR的平均值得到。使用四个尾部切片，我们使用三个VaR分位数(n-1)：
    
    ES = (3.50 + 3.90 + 4.40 + 5.00) / 4 = 4.20
    
    注意95\%的VaR不包含在计算中，因为ES衡量的是超过95\%VaR的平均损失。当尾部切片数量增加时，平均ES会增加，因为计算中包含了更多更高置信水平的VaR。
    
    \textbf{选项A错误。} 预期损失计算结果不正确，应为4.20。
    
    \textbf{选项B错误。} 预期损失计算结果不正确，且增加切片数量时ES会增加而不是减少。
    
    \textbf{选项D错误。} 增加切片数量时ES会增加而不是减少。
\end{explanation}
\checklist{预期损失计算与尾部切片数量影响}



\begin{question}
    A portfolio manager has estimated a fund's 1-year VaR for various confidence levels as shown in the table below:

    \begin{table}[H]
        \centering
        \begin{tabular}{|c|c|}
            \hline
            Confidence Level & Tail VaR \\
            \hline
            95.5\% & 2.1000 \\
            96.0\% & 2.1800 \\
            96.5\% & 2.2700 \\
            97.0\% & 2.3800 \\
            97.5\% & 2.5000 \\
            98.0\% & 2.6700 \\
            98.5\% & 2.8900 \\
            99.0\% & 3.1800 \\
            99.5\% & 3.6400 \\
            \hline
        \end{tabular}
    \end{table}

    What is the best estimate of the 1-year 97.5\% expected shortfall?
\end{question}

\begin{choices}
    \item 2.5000
    \item 2.6700
    \item 3.0850
    \item 3.3450
\end{choices}

\begin{correctanswer}
    D
\end{correctanswer}

\begin{explanation}
    \textbf{选项D正确。}
    
    97.5\%置信水平下的预期损失通过计算97.5\%以上所有VaR的平均值得到：
    
    ES = (2.6700 + 2.8900 + 3.1800 + 3.6400) / 4 = 3.3450
    
    \textbf{选项A错误。} 这是97.5\%置信水平下的VaR值，不是预期损失。
    
    \textbf{选项B错误。} 这是98.0\%置信水平下的VaR值，不是预期损失。
    
    \textbf{选项C错误。} 预期损失计算结果不正确，应为3.3450。
\end{explanation}
\checklist{预期损失计算：取尾部VaR平均值}


\begin{question}
    Assume the profit/loss distribution for XYZ is normally distributed with an annual mean of \$30 million and a standard deviation of \$15 million. The 5\% VaR is calculated and interpreted as which of the following statements?
\end{question}

\begin{choices}
    \item 5\% probability of losses of at least \$5.25 million
    \item 5\% probability of earnings of at least \$5.25 million
    \item 95\% probability of losses of at least \$5.25 million
    \item 95\% probability of earnings of at least \$5.25 million
\end{choices}

\begin{correctanswer}
    D
\end{correctanswer}

\begin{explanation}
    \textbf{选项D正确。}
    
    95\%置信水平下的VaR计算为：
    
    VaR = -30 million + 1.65 × 15 million = -\$5.25 million
    
    由于预期利润为正且VaR为负，这意味着：
    - 有5\%的概率XYZ的盈利少于+\$5.25 million
    - 这等价于说有95\%的概率XYZ至少盈利+\$5.25 million
    
    \textbf{选项A错误。} 5\%VaR不表示损失概率，而是表示盈利低于某个水平的概率。
    
    \textbf{选项B错误。} 5\%VaR表示盈利少于5.25 million的概率，不是至少盈利的概率。
    
    \textbf{选项C错误。} 95\%VaR不表示损失概率，而是表示盈利至少达到某个水平的概率。
\end{explanation}
\checklist{VaR在正态分布下的概率解释}


\begin{question}
    Which of the following statements is not an advantage of spectral risk measures over expected shortfall? Spectral risk measures:
\end{question}

\begin{choices}
    \item incorporate risk aversion preferences into the risk measurement
    \item are a special case of expected shortfall measures
    \item allow for customization of risk weights based on investor preferences
    \item provide more stable risk estimates through better weighting of observations
\end{choices}

\begin{correctanswer}
    B
\end{correctanswer}

\begin{explanation}
    \textbf{选项B正确。}
    
    这个说法是错误的，因为它颠倒了关系。预期损失实际上是谱风险测度的一个特例，而不是相反。谱风险测度更加一般化，包含了预期损失作为特定情况。
    
    \textbf{选项A错误。} 这是谱风险测度的优势，它可以将风险厌恶偏好纳入风险度量中。
    
    \textbf{选项C错误。} 这是谱风险测度的优势，它允许根据投资者偏好定制风险权重。
    
    \textbf{选项D错误。} 这是谱风险测度的优势，它通过更好的观测值加权提供更稳定的风险估计。
\end{explanation}
\checklist{谱风险测度与预期损失的关系}

\begin{question}
    Suppose an investment fund has a two-asset portfolio with \$12 million in asset A and \$8
    million in asset B. The portfolio correlation is 0.40, and the daily standard deviation of
    returns for asset A and B are 3.0\% and 1.5\%, respectively. What is the 10-day value at risk (VaR)
    of this portfolio at a 99\% confidence level (α= 2.33)?
\end{question}

\begin{choices}
    \item \$2.847 million.
    \item \$3.245 million.
    \item \$5.280 million.
    \item \$6.104 million.
\end{choices}

\begin{correctanswer}
    A
\end{correctanswer}

\begin{explanation}
    \textbf{选项A正确。}
    
    步骤1：计算协方差矩阵
    \[
    \text{cov}_{11} = \sigma_1^2 = 0.030^2 = 0.0009
    \]
    \[
    \text{cov}_{22} = \sigma_2^2 = 0.015^2 = 0.000225
    \]
    \[
    \text{cov}_{12} = \rho_{12} \times \sigma_1 \times \sigma_2 = 0.40 \times 0.030 \times 0.015 = 0.00018
    \]
    
    步骤2：计算投资组合标准差
    \[
    \sigma_P = \sqrt{w_1^2\sigma_1^2 + w_2^2\sigma_2^2 + 2w_1w_2\rho\sigma_1\sigma_2}
    \]
    \[
    \sigma_P = \sqrt{(0.6^2 \times 0.0009) + (0.4^2 \times 0.000225) + (2 \times 0.6 \times 0.4 \times 0.00018)} = 0.0389
    \]
    
    步骤3：计算10天VaR
    \[
    \text{VaR}_P = \sigma_P \alpha \sqrt{n} = 0.0389 \times 2.33 \times \sqrt{10} = \$2.847 \text{ million}
    \]
    
    \textbf{选项B错误。} 计算错误，可能是协方差矩阵计算有误。
    
    \textbf{选项C错误。} 使用了年化VaR而不是10天VaR，忘记乘以√10。
    
    \textbf{选项D错误。} 使用了错误的置信水平或VaR公式计算错误。
\end{explanation}
\checklist{投资组合VaR计算：考虑权重、相关性和波动率}


\begin{question}
    A portfolio analyst is estimating the market risk of a fund using both the arithmetic return with normal distribution assumption and the geometric returns with lognormal distribution assumptions. The analyst gathers the following data on the fund:
    \begin{itemize}
        \item Annualized average of arithmetic returns: 18\%
        \item Annualized standard deviation of arithmetic returns: 40\%
        \item Annualized average of geometric returns: 0.42\%
        \item Annualized standard deviation of geometric returns: 50\%
        \item Current fund value: USD 10,000,000
        \item Trading days in a year: 252
    \end{itemize}
    Assuming both daily arithmetic returns and daily geometric returns are serially independent, which of the following statements is correct?
\end{question}

\begin{choices}
    \item 1-day normal 95\% VaR = 4.85\% and 1-day lognormal 95\% VaR = 3.95\%
    \item 1-day normal 95\% VaR = 3.95\% and 1-day lognormal 95\% VaR = 4.85\%
    \item 1-day normal 95\% VaR = 4.85\% and 1-day lognormal 95\% VaR = 4.89\%
    \item 1-day normal 95\% VaR = 3.95\% and 1-day lognormal 95\% VaR = 3.93\%
\end{choices}

\begin{correctanswer}
    B
\end{correctanswer}

\begin{explanation}
    \textbf{选项B正确。}
    
    步骤1：计算正态分布VaR
    \[
    -\left[\frac{0.18}{252} - 1.645 \times \frac{0.40}{\sqrt{252}}\right] = 3.95\%
    \]
    
    步骤2：计算对数正态分布VaR
    \[
    1 - \exp\left[\frac{0.0042}{252} - \frac{0.50 \times 1.645}{\sqrt{252}}\right] = 4.85\%
    \]
    
    \textbf{选项A错误。} 正态VaR和对数正态VaR的数值颠倒了。
    
    \textbf{选项C错误。} 正态VaR计算错误，应为3.95\%而不是4.85\%。
    
    \textbf{选项D错误。} 对数正态VaR计算错误，应为4.85\%而不是3.93\%。
\end{explanation}
\checklist{正态与对数正态分布VaR计算公式差异}


\begin{question}
    It is not always apparent how risk should be quantified for a given financial institution when there are many
    different possible risk measures to consider. Prior to defining specific measures, one should
    be aware of the general characteristics of ideal risk measures. Such measures should be
    intuitive, stable, easy to understand, coherent, and interpretable in economic terms. In
    addition, the risk decomposition process must be simple and meaningful for a given risk
    measure. Standard deviation, value at risk (VaR), expected shortfall (ES), and spectral and
    distorted risk measures are commonly used measures to calculate economic capital. However, it
    is not easy to select a risk measure to calculate economic capital, as each measure has its
    respective pros and cons. Which of the following statements pertaining to the pros and cons of
    these risk measures is not accurate?
\end{question}

\begin{choices}
    \item Standard deviation does not have the property of monotonicity, and therefore, it is not
    coherent.
    \item VaR does not have the property of subadditivity, and therefore; it is not coherent.
    \item ES is not stable regardless of the loss distribution.
    \item Spectral and distorted risk measures are neither intuitive nor commonly used in practice.
\end{choices}

\begin{correctanswer}
    C
\end{correctanswer}

\begin{explanation}
    \textbf{选项C正确。}
    
    预期损失(ES)的稳定性取决于损失分布，而不是绝对不稳定。在许多情况下，ES比VaR更稳定，特别是在处理厚尾分布时。
    
    \textbf{选项A错误。} 标准差确实缺乏单调性，因此不是一致的，这个说法是正确的。
    
    \textbf{选项B错误。} VaR确实缺乏次可加性，因此不是一致的，这个说法是正确的。
    
    \textbf{选项D错误。} 谱风险测度和扭曲风险测度确实不够直观且在实践中使用较少，这个说法是正确的。
\end{explanation}
\checklist{风险度量特性：ES稳定性取决于分布}


\begin{question}
    When selecting risk measures for economic capital calculation, financial institutions need to consider various characteristics such as intuitiveness, stability, coherence, and economic interpretability. Which of the following statements about the properties of different risk measures is not accurate?
\end{question}

\begin{choices}
    \item Standard deviation lacks monotonicity and is therefore not a coherent risk measure
    \item VaR lacks subadditivity and is therefore not a coherent risk measure
    \item Expected shortfall is always stable regardless of the underlying loss distribution
    \item Spectral and distorted risk measures are less intuitive and less commonly used in practice
\end{choices}

\begin{correctanswer}
    C
\end{correctanswer}

\begin{explanation}
    \textbf{选项C正确。}
    
    预期损失(ES)的稳定性取决于基础损失分布，而不是绝对稳定。ES的稳定性特别受到损失分布尾部行为的影响。
    
    \textbf{选项A错误。} 标准差确实缺乏单调性，因此不是一致的风险度量。单调性要求如果一个投资组合的损失总是大于另一个，其风险度量应该更高。标准差可能违反这一性质。
    
    \textbf{选项B错误。} VaR确实违反次可加性，因此不是一致的风险度量。次可加性要求组合投资组合的风险不应超过个别风险之和。VaR可能显示组合投资组合的风险高于个别风险之和。
    
    \textbf{选项D错误。} 谱风险测度和扭曲风险测度确实不够直观且在实践中使用较少。虽然理论上合理，但比VaR或ES等简单度量更复杂，且由于参数估计困难而在实践中较少使用。
\end{explanation}
\checklist{风险度量特性：ES稳定性取决于分布}


\begin{question}
    Which of the following is a true statement about expected shortfall (ES)?
\end{question}

\begin{choices}
    \item ES is a coherent spectral measure which gives equal weight to the tail quantiles
    \item ES is a coherent spectral measure which gives increasingly greater weight to higher tail quantiles
    \item ES is a coherent spectral measure but gives decreasingly less weight to higher tail quantiles
    \item ES is coherent, VaR is not coherent, and neither are spectral measures
\end{choices}

\begin{correctanswer}
    A
\end{correctanswer}

\begin{explanation}
    \textbf{选项A正确。}
    
    预期损失(ES)是一个一致的谱风险测度，它对所有尾部分位数赋予相等的权重。这是ES区别于其他谱风险测度的关键特征。等权重特性使ES比其他风险度量更稳定且更容易解释。
    
    \textbf{选项B错误。} ES并不对更高的尾部分位数赋予递增的更大权重。实际上，ES对所有尾部分位数赋予相等权重，这是其定义特征之一。
    
    \textbf{选项C错误。} ES并不对更高的尾部分位数赋予递减的更小权重。ES在所有尾部分位数中保持等权重，这对保持其一致性性质至关重要。
    
    \textbf{选项D错误。} 这个说法错误地描述了ES、VaR和谱风险测度之间的关系。虽然ES是一致的而VaR不是，但谱风险测度在特定条件下可以是一致的。ES实际上是谱风险测度的一个特例。
\end{explanation}
\checklist{ES作为谱风险测度的等权重特性}



\begin{question}
    A hedge fund uses VaR to estimate the probability of losses. However, management is concerned that the estimation process gives equal weight to all observations. Therefore, alternative approaches are considered. Which of the following statements accurately describes an alternative weighted historic simulation approach? 
    
    I. Individual observations can be weighted based on volatility by substituting historical returns with volatility-adjusted returns.
    
    II. Historical returns can be revised using correlation-adjusted returns.
\end{question}

\begin{choices}
    \item I only
    \item II only
    \item Both I and II
    \item Neither I nor II
\end{choices}

\begin{correctanswer}
    C
\end{correctanswer}

\begin{explanation}
    \textbf{选项C正确。}
    
    两个陈述都是正确的：
    
    陈述I正确描述了波动率加权历史模拟法，其中收益率根据当前波动率水平进行调整。
    
    陈述II正确描述了相关性加权历史模拟法，其中收益率进行调整以反映当前的相关性结构。
    
    \textbf{选项A错误。} 只承认波动率加权方法但忽略了相关性加权方法。两种方法都是等权重历史模拟法的有效替代方案。
    
    \textbf{选项B错误。} 只承认相关性加权方法但忽略了波动率加权方法。两种方法都是等权重历史模拟法的有效替代方案。
    
    \textbf{选项D错误。} 错误地拒绝了两种有效的加权历史模拟法方法。波动率加权和相关性加权都是改进VaR估计的既定方法。
\end{explanation}
\checklist{加权历史模拟法的波动率和相关性调整方法}


\begin{question}
    The accuracy of a value at risk (VaR) measure:
\end{question}

\begin{choices}
    \item Is included in the statistic.
    \item Is complete because the process is deterministic.
    \item Is one minus the probability level.
    \item Can only be ascertained after the fact.
\end{choices}

\begin{correctanswer}
    D
\end{correctanswer}

\begin{explanation}
    \textbf{选项D正确。}
    
    VaR的准确性只能通过观察实际损失后验证，需要回测来验证其可靠性。这是VaR的一个关键局限性。
    
    \textbf{选项A错误。} VaR统计量本身不包含准确性度量，准确性需要单独验证。
    
    \textbf{选项B错误。} VaR计算不是确定性的，涉及假设和估计。
    
    \textbf{选项C错误。} 一减去概率水平是置信水平，不是准确性度量。
\end{explanation}
\checklist{VaR局限性：准确性只能事后验证}



\begin{question}
    Michael Wang is evaluating two investment funds for an investment of \$250,000. Fund A has
    \$40,000,000 in assets, an annual expected return of 18 percent, and an annual standard
    deviation of 25 percent. Fund B has \$15,000,000 in assets, an annual expected return of
    16 percent, and an annual standard deviation of 20 percent. What is the daily value at risk
    (VaR) of Wang's portfolio at a 5 percent probability if he invests his money in fund A?
\end{question}

\begin{choices}
    \item \$4,125.
    \item \$33,000.
    \item \$4,750.
    \item \$57,000.
\end{choices}

\begin{correctanswer}
    C
\end{correctanswer}

\begin{explanation}
    \textbf{选项C正确。}
    
    步骤1：计算日度波动率
    \[
    \sigma = \frac{25\%}{\sqrt{250}} = 0.0158
    \]
    
    步骤2：计算日度收益率
    \[
    \mu = \frac{18\%}{250} = 0.00072
    \]
    
    步骤3：计算日度VaR
    \[
    \text{VaR} = 250{,}000 \times \left| 0.00072 - 1.645 \times 0.0158 \right| = 4{,}750
    \]
    
    \textbf{选项A错误。} 日度波动率计算错误或使用了错误的缩放因子。
    
    \textbf{选项B错误。} 使用了年化VaR而不是日度VaR，忘记按交易日数缩放。
    
    \textbf{选项D错误。} 使用了错误的置信水平或VaR公式计算错误。
\end{explanation}
\checklist{VaR计算：年化转日度需除以根号250}


\begin{question}
    Statement 1: Differences in the use of confidence intervals and time horizon can cause significant variability in VaR estimates as there is lack of uniformity in practice. 
    
    Statement 2: Standardization of confidence interval and time horizon would eliminate most of the variability in VaR estimates. 
    
    The statements are most likely correct with regard to:
   
\end{question}

\begin{choices}
    \item Statement 1 only.
    \item Statement 2 only.
    \item Both statements.
    \item Neither statement.
\end{choices}

\begin{correctanswer}
    A
\end{correctanswer}

\begin{explanation}
    \textbf{选项A正确。}
    
    陈述1正确，因为风险度量的变异性，包括置信水平和时间跨度使用缺乏一致性，可能导致VaR估计的变异性。
    
    陈述2错误，因为其他因素也可能导致变异性，包括分析时间序列的长度、矩估计方法、映射技术、衰减因子和模拟次数等。
    
    \textbf{选项B错误。} 陈述2不正确，标准化置信水平和时间跨度不会消除大部分变异性。
    
    \textbf{选项C错误。} 陈述2不正确，因此两个陈述不都正确。
    
    \textbf{选项D错误。} 陈述1正确，因此两个陈述不都错误。
\end{explanation}
\checklist{VaR估计变异性：置信水平和时间跨度的影响}


\begin{question}
    When selecting risk measures for economic capital calculation, financial institutions face multiple options, each with its own advantages and limitations. Which of the following statements about the properties of different risk measures is not accurate?
\end{question}

\begin{choices}
    \item Standard deviation lacks monotonicity and is therefore not a coherent risk measure
    \item VaR lacks subadditivity and is therefore not a coherent risk measure
    \item Expected shortfall is always stable regardless of the underlying loss distribution
    \item Spectral and distorted risk measures are less intuitive and less commonly used in practice
\end{choices}

\begin{correctanswer}
    C
\end{correctanswer}

\begin{explanation}
    \textbf{选项C正确。}
    
    预期损失(ES)的稳定性取决于基础损失分布，而不是绝对稳定。ES的稳定性特别受到损失分布尾部行为和估计中使用的数据质量的影响。
    
    \textbf{选项A错误。} 标准差确实违反一致风险度量的单调性。单调性要求如果一个投资组合的损失总是大于另一个，其风险度量应该更高。标准差可能违反这一性质，因为它衡量的是离散程度而不是下行风险。
    
    \textbf{选项B错误。} VaR确实违反一致风险度量的次可加性。次可加性要求组合投资组合的风险不应超过个别风险之和。VaR可能显示组合投资组合的风险高于个别风险之和，这对风险管理来说是反直觉的。
    
    \textbf{选项D错误。} 谱风险测度和扭曲风险测度确实不够直观且在实践中使用较少。虽然理论上合理，但比VaR或ES等简单度量更复杂，且由于参数估计困难和向利益相关者传达结果的挑战而在实践中较少使用。
\end{explanation}
\checklist{风险度量特性：ES稳定性取决于分布}


\begin{question}
    A fund manager is calculating the lognormal VaR for a USD 80 million equity portfolio. The portfolio's log returns follow a normal distribution with an annualized mean of 15\% and an annualized standard deviation of 25\%. Assuming 252 trading days in a year, which of the following is closest to the correct estimate for the 1-day 95\% lognormal VaR?
\end{question}

\begin{choices}
    \item USD 1,850,000
    \item USD 2,000,000
    \item USD 2,150,000
    \item USD 2,300,000
\end{choices}

\begin{correctanswer}
    B
\end{correctanswer}

\begin{explanation}
    \textbf{选项B正确。}
    
    对数正态VaR计算为：
    \[
    \text{VaR} = 80,000,000 \times (1 - e^{0.15/252 - 1.645 \times 0.25/\sqrt{252}}) = 2,000,000
    \]
    
    这个计算正确地考虑了收益率的对数正态分布和参数的日度缩放。
    
    \textbf{选项A错误。} 低估了VaR。1,850,000的结果过低，可能是由于对数正态VaR公式应用错误。
    
    \textbf{选项C错误。} 高估了VaR。2,150,000的计算过高，可能是由于参数缩放错误或对数正态变换应用不当。
    
    \textbf{选项D错误。} 显著高估了VaR。2,300,000的计算过高，表明参数缩放或对数正态变换过程中可能存在错误。
\end{explanation}
\checklist{对数正态分布VaR计算：参数年化转日化}






\newpage



\section{考点2:非参数法}
\setcounter{question}{0}
\begin{question}
    Which of the following approaches is not a valid improvement to the traditional historical simulation method for VaR estimation?
\end{question}

\begin{choices}
    \item Volatility-weighted historical simulation
    \item Age-weighted historical simulation
    \item Market-weighted historical simulation
    \item Correlation-weighted historical simulation
\end{choices}

\begin{correctanswer}
    C
\end{correctanswer}

\begin{explanation}
    \textbf{选项C正确。}
    
    "市场加权历史模拟法"不是历史模拟法的标准或有效改进方法。虽然波动率加权、时间加权和相关性加权是既定方法，但在风险管理实践中没有广泛接受的"市场加权"历史模拟法。
    
    \textbf{选项A错误。} 波动率加权历史模拟法是有效的改进方法。它根据当前波动率水平调整历史收益率，使VaR估计对变化的市场条件更敏感。这种方法有助于解决金融市场中波动率聚集的问题。
    
    \textbf{选项B错误。} 时间加权历史模拟法是有效的改进方法。它为更近期的观测值分配更高权重，反映近期市场条件对未来风险估计更相关的信念。这种方法有助于解决市场制度随时间变化的问题。
    
    \textbf{选项D错误。} 相关性加权历史模拟法是有效的改进方法。它调整历史收益率以反映资产间的当前相关性结构，使VaR估计在捕捉当前市场依赖性方面更准确。这种方法有助于解决金融市场中相关性模式变化的问题。
\end{explanation}
\checklist{历史模拟法改进方法：波动率、时间、相关性加权}


\begin{question}
    What is the primary purpose of a quantile-quantile (Q-Q) plot in risk management?
\end{question}

\begin{choices}
    \item To test whether an empirical distribution matches a theoretical distribution
    \item To test whether a theoretical distribution matches an empirical distribution
    \item To identify which theoretical distribution best fits the empirical data
    \item To identify which empirical distribution best fits the theoretical model
\end{choices}

\begin{correctanswer}
    C
\end{correctanswer}

\begin{explanation}
    \textbf{选项C正确。}
    
    Q-Q图主要用于识别哪个理论分布最适合经验数据。该图将经验分布的分位数与理论分布的分位数进行比较。如果点形成一条直线，表明理论分布对经验数据拟合良好。这在风险管理中选择适当的分布假设特别有用。
    
    \textbf{选项A错误。} Q-Q图不用于正式统计检验。它们主要是用于分布比较和识别的可视化工具，而不是假设检验。
    
    \textbf{选项B错误。} Q-Q图不是为检验理论分布而设计的。它们用于视觉评估理论分布对经验数据的拟合程度。
    
    \textbf{选项D错误。} Q-Q图不用于识别经验分布。它们用于将经验数据与理论分布进行比较，以找到最佳的理论拟合。
\end{explanation}
\checklist{Q-Q图用途：识别最适合经验数据的理论分布}





\begin{question}
    Which of the following statements correctly describes the volatility-weighting approach in VaR calculation?
\end{question}

\begin{choices}
    \item Historical returns are adjusted, and the VaR calculation becomes more complex
    \item Historical returns are adjusted, but the VaR calculation method remains unchanged
    \item Current returns are adjusted, and the VaR calculation becomes more complex
    \item Current returns are adjusted, but the VaR calculation method remains unchanged
\end{choices}

\begin{correctanswer}
    B
\end{correctanswer}

\begin{explanation}
    \textbf{选项B正确。}
    
    波动率加权方法只通过将历史收益率乘以当前波动率与历史波动率的比率来调整历史收益率。公式为：
    \[
    \text{调整后收益率}_t = \text{历史收益率}_t \times \frac{\text{当前波动率}}{\text{历史波动率}_t}
    \]
    
    调整后，使用标准历史模拟法计算VaR，计算程序本身没有变化。
    
    \textbf{选项A错误。} 虽然历史收益率确实在波动率加权中被调整，但VaR计算程序本身不会变得更复杂。调整应用于输入数据，而不是计算方法。
    
    \textbf{选项C错误。} 波动率加权调整的是历史收益率，而不是当前收益率。目的是使历史数据与当前市场条件更相关。
    
    \textbf{选项D错误。} 波动率加权不调整当前收益率。它调整历史收益率以反映当前波动率水平，同时保持相同的VaR计算方法。
\end{explanation}
\checklist{波动率加权：调整历史收益率但保持VaR计算程序不变}

\begin{question}
    The following table shows the six lowest returns for a portfolio with their hybrid weights:

    \begin{table}[H]
        \centering
        \begin{tabular}{|c|c|c|c|}
            \hline
            \multicolumn{4}{|c|}{Six lowest returns with hybrid weightings} \\
            \hline
            & Six lowest returns & Hybrid weight & Hybrid Cumulative weight \\
            \hline
            1 & -5.20\% & 0.0160 & 0.0160 \\
            2 & -4.80\% & 0.0150 & 0.0310 \\
            3 & -4.50\% & 0.0100 & 0.0410 \\
            4 & -4.20\% & 0.0120 & 0.0530 \\
            5 & -4.00\% & 0.0080 & 0.0610 \\
            6 & -3.80\% & 0.0040 & 0.0650 \\
            \hline
        \end{tabular}
        \caption{Six lowest returns with hybrid weightings}
    \end{table}

    What is the 5\% VaR under the hybrid approach?
\end{question}

\begin{choices}
    \item -4.00\%
    \item -4.15\%
    \item -4.30\%
    \item -4.20\%
\end{choices}

\begin{correctanswer}
    C
\end{correctanswer}

\begin{explanation}
    \textbf{选项C正确。}
    
    第四个收益率(-4.20\%)的累积权重为5.30\%
    第五个收益率(-4.00\%)的累积权重为6.10\%
    5\%VaR落在这两个点之间
    
    使用线性插值：
    \[
    \text{VaR} = -4.20\% + \frac{5\% - 5.30\%}{6.10\% - 5.30\%} \times (-4.00\% - (-4.20\%)) = -4.30\%
    \]
    
    \textbf{选项A错误。} -4.00\%对应第五个最低收益率，累积权重为6.10\%，高于5\%阈值。5\%VaR应该在第四和第五个最低收益率之间。
    
    \textbf{选项B错误。} -4.15\%不是正确的插值结果。插值应该在第四个收益率(-4.20\%, 5.30\%)和第五个收益率(-4.00\%, 6.10\%)之间。
    
    \textbf{选项D错误。} -4.20\%对应第四个最低收益率，累积权重为5.30\%，高于5\%阈值。5\%VaR应该在第四和第五个最低收益率之间。
\end{explanation}
\checklist{混合加权法VaR计算：线性插值法}


\begin{question}
    Which of the following statements about age-weighting in VaR calculation is most accurate?
\end{question}

\begin{choices}
    \item The age-weighting method requires GARCH model estimates for implementation
    \item If the decay factor is close to 1, the data shows strong persistence
    \item The weight for day i is calculated as:
    \[
    W_{(i)} = \frac{\lambda^{i-1} \times (1 - \lambda)}{1 - \lambda^{n}}
    \]
    \item The minimum required number of observations is 250
\end{choices}

\begin{correctanswer}
    B
\end{correctanswer}

\begin{explanation}
    \textbf{选项B正确。}
    
    当衰减因子(λ)接近1时，权重衰减非常缓慢，这意味着较旧的观测值保持显著影响。数据中的持久性反映在权重的缓慢衰减中，这在市场条件相对稳定时特别有用。
    
    \textbf{选项A错误。} 时间加权法不需要GARCH模型估计。它是一种更简单的方法，基于观测值的年龄分配权重，使用衰减因子给近期数据更多权重。
    
    \textbf{选项C错误。} 显示的公式不正确。正确的公式应该在分母中使用n(总观测数)而不是i。正确公式为：
    \[
    W_{(i)} = \frac{\lambda^{i-1} \times (1 - \lambda)}{1 - \lambda^{n}}
    \]
    
    \textbf{选项D错误。} 时间加权法没有250个观测值的严格最低要求。适当的样本量取决于具体应用和市场条件。
\end{explanation}
\checklist{时间加权法：衰减因子对数据持久性的影响}


\begin{question}
    What is the key difference between historical simulation and bootstrapped historical simulation in VaR estimation?
\end{question}

\begin{choices}
    \item Bootstrapping is a non-parametric approach
    \item Bootstrapping can be used to estimate both VaR and ES
    \item Bootstrapping does not require a variance-covariance matrix
    \item Bootstrapping generates multiple samples through resampling
\end{choices}

\begin{correctanswer}
    D
\end{correctanswer}

\begin{explanation}
    \textbf{选项D正确。}
    
    历史模拟法使用单一的历史数据样本，而自举历史模拟法通过对原始数据进行有放回抽样生成多个样本。每个重采样数据集用于计算VaR估计，最终的VaR估计通常是所有重采样VaR的平均值。
    
    \textbf{选项A错误。} 历史模拟法和自举历史模拟法都是非参数方法，这不是两种方法之间的区别特征。
    
    \textbf{选项B错误。} 两种方法都可以用于估计VaR和ES，这是两种方法的共同特征，而不是它们之间的差异。
    
    \textbf{选项C错误。} 两种方法都不需要方差-协方差矩阵，两种方法都直接使用历史收益率数据。
\end{explanation}
\checklist{自举历史模拟法：通过重采样生成多个样本}

\begin{question}
    Michael Wang has collected 1,500 daily observations of equity returns. He is concerned about the skewness in the data and decides to use historical simulation with bootstrapping to estimate the 5\% VaR. Which of the following steps is most likely to be part of the estimation procedure?
\end{question}

\begin{choices}
    \item Filter the data to remove outliers
    \item Perform repeated sampling with replacement
    \item Identify the tail region by reordering the original data
    \item Apply a weighting scheme to reduce the impact of older data
\end{choices}

\begin{correctanswer}
    B
\end{correctanswer}

\begin{explanation}
    \textbf{选项B正确。}
    
    自举法涉及对原始数据集进行有放回重复抽样。每个重采样数据集用于计算5\%VaR估计，最终的VaR估计是所有重采样VaR的平均值。这个过程有助于在不做分布假设的情况下处理数据中的偏度。
    
    \textbf{选项A错误。} 自举法不涉及过滤或移除异常值。该方法使用完整数据集，包括所有观测值，以保持收益率的真实分布。
    
    \textbf{选项C错误。} 自举法不需要重新排序原始数据。重采样过程是随机的，可以从数据集中选择任何观测值。
    
    \textbf{选项D错误。} 自举法不涉及对观测值加权。每个观测值在每次重采样中被选中的概率相等。
\end{explanation}
\checklist{自举历史模拟法：有放回重复抽样计算平均VaR}


\begin{question}
    A portfolio manager is comparing parametric and non-parametric approaches for VaR calculation and is concerned about the characteristics of the loss data. Which of the following distribution characteristics would make parametric approaches the preferred method?
\end{question}

\begin{choices}
    \item Skewness in the distribution
    \item Fat tails in the distribution
    \item Scarcity of high magnitude loss events
    \item Heteroskedasticity in the distribution
\end{choices}

\begin{correctanswer}
    C
\end{correctanswer}

\begin{explanation}
    \textbf{选项C正确。}
    
    当数据中高幅度损失事件较少时，非参数方法可能低估风险。参数方法可以通过做分布假设更好地外推风险。极端事件的稀缺性使得非参数方法难以准确估计尾部风险，而参数方法在这种情况下可以提供更稳定的风险估计。
    
    \textbf{选项A错误。} 分布中的偏度实际上是选择非参数方法的理由。参数方法通常假设正态分布，无法充分捕捉数据中的偏度。
    
    \textbf{选项B错误。} 分布中的厚尾更好地由非参数方法处理。参数方法由于其正态分布假设，在存在厚尾时通常低估风险。
    
    \textbf{选项D错误。} 异方差性(波动率变化)更好地由非参数方法处理。参数方法通常假设恒定波动率，当波动率随时间变化时可能导致不准确的风险估计。
\end{explanation}
\checklist{参数法优势：极端损失事件较少时的适用性}


\begin{question}
    Which of the following statements is true regarding the bootstrap simulation method used in VaR estimation? 
    
    Ⅰ. Bootstrapping uses actual market data 
    
    Ⅱ. The bootstrapping method always uses a time horizon based on the time scale of the historical data 
    
    Ⅲ. Bootstrapping is based on synthesis of normally distributed random numbers
\end{question}

\begin{choices}
    \item I only
    \item II only
    \item I and III
    \item I, II and III
\end{choices}

\begin{correctanswer}
    A
\end{correctanswer}

\begin{explanation}
    \textbf{选项A正确。}
    
    陈述I正确：自举法使用实际历史市场数据。
    
    陈述II错误：自举法可以使用与历史数据不同的时间尺度，使其比陈述II暗示的更灵活。
    
    陈述III错误：自举法使用实际市场数据，而不是正态分布随机数。使用正态分布随机数是蒙特卡洛模拟的特征，而不是自举法。
    
    \textbf{选项B错误。} 陈述II错误，自举法可以使用与历史数据不同的时间尺度。
    
    \textbf{选项C错误。} 陈述III错误，自举法使用实际市场数据，而不是正态分布随机数。
    
    \textbf{选项D错误。} 陈述II和III都错误，只有陈述I关于使用实际市场数据是正确的。
\end{explanation}
\checklist{自举法特点：使用实际市场数据重抽样}


\begin{question}
    Robert Liu has collected a large dataset of daily market returns for four emerging markets. He is concerned about the non-normal skew in the data and is considering non-parametric estimation methods. Liu discusses the procedure with his colleague Jennifer Zhang. During their discussion, Zhang makes the following statements:

    I. Age-weighted historical simulation reduces the impact of older observations only after surpassing a user-defined threshold.
    
    II. Volatility-weighted historical simulation adjusts historic returns with an additive volatility adjustment.
    
    III. Filtered historical estimation combines sophisticated parametric and dynamic volatility estimation techniques.

    How many of Ms. Zhang's statements are correct?
\end{question}

\begin{choices}
    \item Zero
    \item One
    \item Two
    \item Three
\end{choices}

\begin{correctanswer}
    A
\end{correctanswer}

\begin{explanation}
    \textbf{选项A正确。}
    
    所有三个陈述都是错误的：
    
    陈述I错误：时间加权历史模拟法通过恒定衰减因子减少每个观测值的权重，而不是通过阈值。
    
    陈述II错误：波动率加权历史模拟法使用乘数调整(将收益率乘以当前波动率与历史波动率的比率)，而不是加法调整。
    
    陈述III错误：过滤历史模拟法将历史模拟法与条件波动率模型结合，而不是与复杂的参数技术结合。
    
    \textbf{选项B错误。} 没有陈述是正确的，每个陈述都包含对各自方法的基本误解。
    
    \textbf{选项C错误。} 没有陈述是正确的，这些陈述错误地描述了每种方法的核心特征。
    
    \textbf{选项D错误。} 没有陈述是正确的，所有三个陈述在描述方法时都包含重大错误。
\end{explanation}
\checklist{历史模拟法改进方法：时间加权、波动率加权、过滤历史模拟}


\begin{question}
    Which of the following statements is not an advantage of non-parametric simulation methods in risk measurement?
\end{question}

\begin{choices}
    \item Data that requires adjustments is often readily available
    \item Intuitive and often computationally simple
    \item Not hindered by parametric violations of skewness
    \item Can accommodate more complex analysis
\end{choices}

\begin{correctanswer}
    A
\end{correctanswer}

\begin{explanation}
    \textbf{选项A正确。}
    
    这不是非参数方法的优势。实际上，非参数方法的一个关键优势是它们使用不需要调整的原始数据。该陈述错误地暗示需要调整的数据是一个优势，而实际上这是一个局限性。
    
    \textbf{选项B错误。} 这是非参数方法的真正优势。这些方法通常更容易理解和实施，与参数方法相比通常需要更少的计算复杂性。
    
    \textbf{选项C错误。} 这是非参数方法的真正优势。这些方法可以在不做分布假设的情况下处理偏斜分布，这比通常假设正态分布的参数方法有显著优势。
    
    \textbf{选项D错误。} 这是非参数方法的真正优势。这些方法可以应用于复杂的金融工具和投资组合，而不需要特定的分布假设。
\end{explanation}
\checklist{非参数方法优势：使用原始数据、计算简单、处理偏斜分布}


\begin{question}
    Crude oil prices show seasonal volatility patterns, with higher volatility during hurricane season. Which of the following statements is correct if VaR is estimated using unweighted historical simulation with a three-year sample period?
\end{question}

\begin{choices}
    \item We will overstate VaR in the off-season and understate VaR during hurricane season
    \item We will overstate VaR in the off-season and overstate VaR during hurricane season
    \item We will understate VaR in the off-season and understate VaR during hurricane season
    \item We will understate VaR in the off-season and overstate VaR during hurricane season
\end{choices}

\begin{correctanswer}
    A
\end{correctanswer}

\begin{explanation}
    \textbf{选项A正确。}
    
    未加权历史模拟法计算整个三年期间的平均波动率。这个平均波动率高于实际的淡季波动率(导致高估)，低于实际的飓风季节波动率(导致低估)。该方法无法考虑波动率的季节性模式。
    
    \textbf{选项B错误。} 该选项暗示VaR在两个季节都会被高估。实际上，平均波动率会低于飓风季节波动率，导致飓风季节的低估。
    
    \textbf{选项C错误。} 该选项暗示VaR在两个季节都会被低估。实际上，平均波动率会高于淡季波动率，导致淡季的高估。
    
    \textbf{选项D错误。} 该选项暗示VaR在淡季会被低估，在飓风季节会被高估。实际上，由于未加权历史模拟法的平均效应，情况正好相反。
\end{explanation}
\checklist{未加权历史模拟法局限性：季节性波动处理不当}


\begin{question}
    The historical simulation (HS) approach is based on empirical distributions and can handle multiple risk factors. The RiskMetrics approach assumes normal distributions and uses mapping to equity indices. For which of the following portfolios would the HS approach be more likely to provide an accurate VaR estimate than the RiskMetrics approach?
\end{question}

\begin{choices}
    \item A small number of emerging market securities
    \item A small number of broad market indexes
    \item A large number of emerging market securities
    \item A large number of broad market indexes
\end{choices}

\begin{correctanswer}
    A
\end{correctanswer}

\begin{explanation}
    \textbf{选项A正确。}
    
    新兴市场证券通常具有非正态分布，具有厚尾和偏度特征。小型投资组合不太可能从分散化效应中受益。RiskMetrics的映射方法对非正态分布效果较差。历史模拟法可以在不做分布假设的情况下更好地捕捉收益率的实际分布。
    
    \textbf{选项B错误。} 广泛市场指数倾向于具有更接近正态的分布，使RiskMetrics的正态分布假设更合适。映射方法对这些证券效果良好。
    
    \textbf{选项C错误。} 大量新兴市场证券会从分散化效应中受益，使正态分布假设更合理。在这种情况下RiskMetrics表现会更好。
    
    \textbf{选项D错误。} 大量广泛市场指数投资组合的收益率更接近正态分布，使RiskMetrics的假设更合适。映射方法对这些证券效果良好。
\end{explanation}
\checklist{历史模拟法优势：处理非正态分布小型新兴市场投资组合}

\begin{question}
    Sarah Johnson has collected a large dataset of daily market returns for four emerging markets. She is concerned about the non-normal skew in the data and is considering non-parametric estimation methods. Johnson discusses the procedure with her colleague Michael Brown. During their discussion, Brown makes the following statements:

    I. Age-weighted historical simulation reduces the impact of older observations only after surpassing a user-defined threshold.
    
    II. Volatility-weighted historical simulation adjusts historic returns with an additive volatility adjustment.
    
    III. Filtered historical estimation combines sophisticated parametric and dynamic volatility estimation techniques.

    How many of Mr. Brown's statements are correct?
\end{question}

\begin{choices}
    \item Zero
    \item One
    \item Two
    \item Three
\end{choices}

\begin{correctanswer}
    A
\end{correctanswer}

\begin{explanation}
    \textbf{选项A正确。}
    
    所有三个陈述都是错误的：
    
    陈述I错误：时间加权历史模拟法使用恒定衰减因子随时间减少每个观测值的权重，而不是基于阈值的方法。
    
    陈述II错误：波动率加权历史模拟法使用乘数调整(将收益率乘以当前波动率与历史波动率的比率)，而不是加法调整。
    
    陈述III错误：过滤历史模拟法将历史模拟法与条件波动率模型结合，而不是与复杂的参数技术结合。
    
    \textbf{选项B错误。} 没有陈述是正确的，每个陈述都包含对各自方法的基本误解。
    
    \textbf{选项C错误。} 没有陈述是正确的，这些陈述错误地描述了每种方法的核心特征。
    
    \textbf{选项D错误。} 没有陈述是正确的，所有三个陈述在描述方法时都包含重大错误。
\end{explanation}
\checklist{历史模拟法改进方法：时间加权、波动率加权、过滤历史模拟}


\begin{question}
    A portfolio manager at a large investment fund wants to improve the accuracy of the fund's ES estimates. The manager has collected the profit and loss data on one of the fund's portfolios over the past 1,200 trading days and plans to apply the bootstrapped historical simulation approach to generate ES estimates. Which of the following statements about the bootstrapped historical approach is correct?
\end{question}

\begin{choices}
    \item A major cost incurred when applying the bootstrapped historical simulation approach is the increased computational time and complex programming required
    \item The non-parametric nature of the bootstrapped historical simulation approach makes this approach more difficult to apply when there are a large number of assets in the portfolio
    \item The number of resamples taken from the data set should be chosen to minimize the bias of the ES estimate
    \item Resamples are obtained by choosing profit and loss observations from the original data set at random with replacement
\end{choices}

\begin{correctanswer}
    D
\end{correctanswer}

\begin{explanation}
    \textbf{选项D正确。}
    
    自举法涉及对原始数据集进行有放回随机抽样。每个观测值在重采样中可以被多次选择。这个过程保持了原始数据的经验分布，重采样过程是自举法的根本。
    
    \textbf{选项A错误。} 自举历史模拟法相对简单易实施，不需要复杂的编程。计算要求适中且可管理。
    
    \textbf{选项B错误。} 自举法的非参数性质实际上使其更容易应用于大型投资组合，因为它不需要复杂的分布假设或协方差矩阵。
    
    \textbf{选项C错误。} 重采样数量通常基于计算资源和期望精度选择，而不是专门为了最小化偏差。ES估计的偏差更多与原始数据的质量相关。
\end{explanation}
\checklist{自举历史模拟法：有放回随机抽样生成样本}

\begin{question}
    A senior analyst at a hedge fund is presenting to a group of junior analysts on the historical simulation (HS) approach for estimating ES. The analyst discusses the assumptions, data requirements, advantages, and disadvantages of the HS approach. Which of the following statements would the analyst be correct to make about an advantage of using the HS approach for estimating ES?
\end{question}

\begin{choices}
    \item Major shifts in financial markets or the factors that influence them can be ignored since any effects of these shifts would already have been included in historical datasets
    \item Any type of investment position can be easily accommodated, from simple linear positions to complex non-linear derivatives positions
    \item ES can be easily and accurately estimated for any holding period up to the full time period of the historical data set
    \item Features of profit and loss distributions such as fat tails and skewness are easily accommodated by making appropriate distributional assumptions
\end{choices}

\begin{correctanswer}
    B
\end{correctanswer}

\begin{explanation}
    \textbf{选项B正确。}
    
    历史模拟法可以在不做分布假设的情况下处理任何类型的头寸。它对线性头寸(股票、债券)和非线性头寸(期权、衍生品)都同样有效。该方法使用实际历史价格变化，使其适用于复杂工具。
    
    \textbf{选项A错误。} 重大市场变化不能被忽略。历史模拟法假设未来收益率将遵循与过去收益率相同的分布，这在重大市场变化期间可能不成立。
    
    \textbf{选项C错误。} 历史模拟法不能轻易估计超过历史数据期间的持有期ES。该方法需要足够的历史数据来支持期望的持有期。
    
    \textbf{选项D错误。} 历史模拟法不做分布假设。它使用收益率的经验分布，自然地捕捉厚尾和偏度等特征，而不需要明确的假设。
\end{explanation}
\checklist{历史模拟法优势：适用于各种投资头寸无需分布假设}



\newpage


\section{考点3:验证VaR模型}
\setcounter{question}{0}

\begin{question}
    Michael Wang, FRM, is testing the capital adequacy of two large investment funds. He has collected the following information on historical returns and standard errors (which are normally distributed). Following current regulations, he will test each fund's VaR at a 95\% confidence level. Which of the following statements best describes his findings?

    Fund A: Mean P/L 18\%, Standard deviation P/L 15\%, Default frequency 1.5\%
    
    Fund B: Mean P/L 12\%, Standard deviation P/L 10\%, Default frequency 3\%
\end{question}

\begin{choices}
    \item Fund A's VaR is roughly two times greater than Fund B's VaR
    \item Fund B's VaR is roughly two times greater than Fund A's VaR
    \item Fund B's VaR is less than Fund A's VaR which is less than two times Fund B's VaR
    \item Fund A's VaR is less than Fund B's VaR which is less than two times Fund A's VaR
\end{choices}

\begin{correctanswer}
    A
\end{correctanswer}

\begin{explanation}
    \textbf{选项A正确。}
    
    对于95\%置信水平，我们使用1.65作为临界值。
    
    基金A的VaR = -18\% + 1.65 × 15\% = 6.75\%
    基金B的VaR = -12\% + 1.65 × 10\% = 4.5\%
    
    基金A的VaR大约是基金B的VaR的1.5倍。
    
    \textbf{选项B错误。} 基金B的VaR(4.5\%)实际上小于基金A的VaR(6.75\%)。
    
    \textbf{选项C错误。} 基金B的VaR小于基金A的VaR，但比率不是如描述的那样。实际比率大约是1.5倍。
    
    \textbf{选项D错误。} 基金A的VaR大于基金B的VaR，而不是小于它。
\end{explanation}
\checklist{正态分布VaR计算：VaR = -μ + z × σ}



\newpage



\section{考点4:极值理论}
\setcounter{question}{0}
\begin{question}
    In modeling extreme values using the generalized extreme value (GEV) distribution, cases where the tail index parameter is less than zero:
\end{question}

\begin{choices}
    \item Cannot be modeled, and they are usually not of interest in financial risk management
    \item Can be modeled, but they are usually not of interest in financial risk management
    \item Can be modeled, and they are usually of the most interest in financial risk management
    \item Cannot be modeled, but would be of interest in financial risk management if they could be modeled
\end{choices}

\begin{correctanswer}
    B
\end{correctanswer}

\begin{explanation}
    \textbf{选项B正确。}
    
    当ξ < 0时，GEV分布具有有限上界，这表示薄尾分布。这种分布在金融风险建模中通常没有用处。金融收益率通常表现出厚尾(ξ > 0)而不是薄尾。
    
    \textbf{选项A错误。} 尾部参数为负(ξ < 0)的情况可以使用GEV分布建模。该分布对所有尾部参数值都有良好定义。
    
    \textbf{选项C错误。} 薄尾分布(ξ < 0)在金融风险管理中不是主要关注点。厚尾分布(ξ > 0)对建模金融风险更相关。
    
    \textbf{选项D错误。} 尾部参数为负的情况可以建模。GEV分布对所有尾部参数值都有效。
\end{explanation}
\checklist{GEV分布尾部参数：负值表示薄尾分布}

\begin{question}
    Let X be a random variable representing the daily loss of your portfolio. The "peaks over threshold" (POT) approach considers a threshold value, u, of X and the distribution of excess losses over this threshold. Which of the following statements about this application of extreme value theory is correct?
\end{question}

\begin{choices}
    \item To apply the POT approach, the distribution of X must be elliptical and known
    \item If X is normally distributed, the distribution of excess losses requires the estimation of only one parameter, $\beta$, which is a positive scale parameter
    \item To apply the POT approach, one must choose a threshold, u, which is high enough that the number of observations in excess of u is zero
    \item As the threshold, u, increases, the distribution of excess losses over u converges to a generalized Pareto distribution
\end{choices}

\begin{correctanswer}
    D
\end{correctanswer}

\begin{explanation}
    \textbf{选项D正确。}
    
    广义帕累托分布是随着阈值增加时超额损失的极限分布。这是极值理论中的一个基本结果。收敛到这个分布与X的基础分布无关。这个性质使POT方法特别适用于建模极端损失。
    
    \textbf{选项A错误。} POT方法不要求基础分布是椭圆的或已知的。它可以应用于任何分布，包括未知分布。
    
    \textbf{选项B错误。} 超额损失分布需要估计两个参数：尺度参数(β)和形状参数(ξ)，无论X的基础分布如何。
    
    \textbf{选项C错误。} 阈值u必须选择得足够高，以便有足够的观测值超过它。如果没有超过阈值的观测值，该方法无法应用。
\end{explanation}
\checklist{超阈值方法：超额损失分布收敛于广义帕累托分布}


\begin{question}
    According to the Fisher-Tippett theorem, as the sample size n gets large, the distribution of extremes converges to:
\end{question}

\begin{choices}
    \item A normal distribution
    \item A uniform distribution
    \item A generalized Pareto distribution
    \item A generalized extreme value distribution
\end{choices}

\begin{correctanswer}
    D
\end{correctanswer}

\begin{explanation}
    \textbf{选项D正确。}
    
    Fisher-Tippett定理指出极值分布收敛于广义极值(GEV)分布。这种收敛与基础分布无关。GEV分布是块最大值的极限分布。这是极值理论中的一个基本结果。
    
    \textbf{选项A错误。} 正态分布适用于中心极限定理，它描述样本均值的分布，而不是极值。Fisher-Tippett定理专门处理极值的分布。
    
    \textbf{选项B错误。} 均匀分布与极值理论无关。Fisher-Tippett定理在其结论中没有提到均匀分布。
    
    \textbf{选项C错误。} 广义帕累托分布用于超阈值(POT)方法，这是极值理论的另一个分支。它来自Pickands-Balkema-de Haan定理，而不是Fisher-Tippett定理。
\end{explanation}
\checklist{Fisher-Tippett定理：极值分布收敛于广义极值分布}


\begin{question}
    A hedge fund with an active position in commodity futures is using the peaks-over-threshold (POT) methodology for estimating VaR and ES at the 99\% confidence level. The fund's risk managers have set a threshold level to evaluate excess losses. The choice of the threshold, they argue, is suitable and consistent with the finding that 4.50\% of the observations are in excess of the threshold value. The risk managers have concluded that the position's VaR using the POT measure is 5.20\%. The VaR estimate is computed from the following parameters and the managers' empirical analysis is based upon the generalized Pareto distribution assumption for the excess losses.
     
    \begin{table}[H]
        \centering
        \begin{tabular}{|l|c|c|}
        \hline
        \textbf{Parameter} & \textbf{Symbol} & \textbf{Value} \\
        \hline
        Loss threshold & $u$ & 3.5 \\
        Number of observations & $N$ & 800 \\
        Number of observations that exceed threshold & $n$ & 36 \\
        Scale & $\beta$ & 0.85 \\
        Shape (tail index) & $\xi$ & 0.25 \\
        \hline
        \end{tabular}
        \end{table}
    
    Given the VaR value and the parameter assumptions, which of the following is correct?
\end{question}

\begin{choices}
    \item Keeping all other parameters constant, increasing the value of the tail index lowers both the ES and the VaR
    \item Keeping all other parameters constant, increasing the loss threshold level increases both the ES and the VaR
    \item The value of ES is 5.35\%
    \item The value of ES is 6.45\%
\end{choices}

\begin{correctanswer}
    B
\end{correctanswer}

\begin{explanation}
    \textbf{选项B正确。}
    
    提高损失阈值(u)会增加VaR和ES。这从POT公式可以看出：
    \[
    \text{VaR} = u + \left( \frac{\beta}{\xi} \right) \left\{ \left[ \frac{N}{n} (1 - \text{confidence level}) \right]^{-\xi} - 1 \right\}
    \]
    \[
    \text{ES} = \frac{\text{VaR}}{1-\xi} + \frac{\beta - \xi u}{1-\xi}
    \]
    
    两个度量都与阈值水平直接相关。
    
    \textbf{选项A错误。} 提高尾部指数参数(ξ)实际上会增加VaR和ES。这是因为更高的尾部指数表示更重的尾部，导致更高的风险度量。
    
    \textbf{选项C错误。} 实际ES值为6.93\%，使用公式计算：
    \[
    \text{ES} = \frac{5.20}{1-0.25} + \frac{0.85-0.25 \times 3.5}{1-0.25} = 6.93\%
    \]
    
    \textbf{选项D错误。} 实际ES值为6.93\%，不是6.45\%。计算显示ES高于两个建议值。
\end{explanation}
\checklist{超阈值方法：提高阈值会增加VaR和ES}

\begin{question}
    Which of the following about the most common distribution used for peaks-over-threshold is false? 
    
    Ⅰ The distribution requires a threshold, shape, and scaling parameter.
    
    Ⅱ The distribution of these extreme values follows the GEV distribution. 
    
    Ⅲ The distribution produces a curve that dips below the normal distribution prior to the tail and then moves above the normal distribution in a curved shape until it reached the extreme tail. 
    
    Ⅳ The distribution provides a more accurate estimate of the event probabilities in the distribution tail, allowing VaR to be computed at high confidence levels.
\end{question}

\begin{choices}
    \item I and II
    \item II only
    \item III and IV
    \item IV only
\end{choices}

\begin{correctanswer}
    A
\end{correctanswer}

\begin{explanation}
    \textbf{选项A正确。}
    
    陈述I正确：广义帕累托分布(GPD)需要三个参数：阈值(u)、形状(ξ)和尺度(β)。
    
    陈述II错误：POT中极值的分布遵循GPD，而不是GEV分布。GEV用于块最大值方法，而不是POT。
    
    \textbf{选项B错误。} 该选项只识别陈述II为错误，但陈述I也是正确答案的一部分。
    
    \textbf{选项C错误。} 陈述III和IV都是正确的：
    - 陈述III正确描述了GPD曲线的行为
    - 陈述IV正确说明GPD为高置信水平VaR提供更好的尾部估计
    
    \textbf{选项D错误。} 陈述IV是正确的。GPD确实为高置信水平VaR计算提供更准确的尾部估计。
\end{explanation}
\checklist{超阈值方法：使用广义帕累托分布而非广义极值分布}



\begin{question}
    The peaks-over-threshold approach generally requires:
\end{question}

\begin{choices}
    \item More estimated parameters than the GEV approach and shares one parameter with the GEV
    \item Fewer estimated parameters than the GEV approach and shares one parameter with the GEV
    \item More estimated parameters than the GEV approach and does not share any parameters with the GEV approach
    \item Fewer estimated parameters than the GEV approach and does not share any parameters with the GEV approach
\end{choices}

\begin{correctanswer}
    B
\end{correctanswer}

\begin{explanation}
    \textbf{选项B正确。}
    
    超阈值方法比广义极值方法需要更少的参数。两种方法都共享形状参数(ξ)。形状参数ξ至关重要，因为它指示尾部行为：
    - ξ > 0表示厚尾
    - ξ = 0表示指数尾
    - ξ < 0表示薄尾
    
    \textbf{选项A错误。} 超阈值方法实际上比广义极值方法需要更少的参数。POT通常使用三个参数(阈值、尺度和形状)，而GEV使用更多参数。
    
    \textbf{选项C错误。} POT不需要比GEV更多的参数，两种方法都共享形状参数ξ。
    
    \textbf{选项D错误。} 虽然POT确实比GEV需要更少的参数，但两种方法都共享形状参数ξ。
\end{explanation}
\checklist{超阈值方法：参数更少但共享形状参数}


\begin{question}
    In setting the threshold in the POT approach, which of the following statements is the most accurate? Setting the threshold relatively high makes the model:
\end{question}

\begin{choices}
    \item More applicable but decreases the number of observations in the modeling procedure
    \item Less applicable and decreases the number of observations in the modeling procedure
    \item More applicable but increases the number of observations in the modeling procedure
    \item Less applicable but increases the number of observations in the modeling procedure
\end{choices}

\begin{correctanswer}
    A
\end{correctanswer}

\begin{explanation}
    \textbf{选项A正确。}
    
    设定更高的阈值使模型更符合极值理论。更高的阈值确保超过它的观测值遵循广义帕累托分布。然而，更高的阈值减少了可用于参数估计的观测值数量。这在理论有效性和统计可靠性之间创造了权衡。
    
    \textbf{选项B错误。} 更高的阈值实际上使模型更符合极值理论，而不是更不适用。
    
    \textbf{选项C错误。} 更高的阈值减少而不是增加可用于建模的观测值数量。
    
    \textbf{选项D错误。} 更高的阈值使模型更符合极值理论，并且减少而不是增加观测值数量。
\end{explanation}
\checklist{超阈值方法阈值设定：提高阈值更符合理论但减少观测值}

\begin{question}
    Given a VaR equal to 3.20, a threshold of 1.5\%, a shape parameter equal to 0.25, and a scale parameter equal to 0.40, what is the expected shortfall?
\end{question}

\begin{choices}
    \item 4.267
    \item 4.533
    \item 4.800
    \item 5.067
\end{choices}

\begin{correctanswer}
    A
\end{correctanswer}

\begin{explanation}
    \textbf{选项A正确。}
    
    使用超阈值方法的预期损失公式：
    \[
    \text{ES} = \frac{\text{VaR}}{1-\xi} + \frac{\beta - \xi u}{1-\xi}
    \]
    
    代入给定值：
    \[
    \text{ES} = \frac{3.20}{1-0.25} + \frac{0.40 - 0.25 \times 1.5}{1-0.25} = 4.267
    \]
    
    其中：VaR = 3.20，ξ(形状参数) = 0.25，β(尺度参数) = 0.40，u(阈值) = 1.5%
    
    \textbf{选项B错误。} 计算结果不正确，应为4.267。
    
    \textbf{选项C错误。} 计算结果不正确，应为4.267。
    
    \textbf{选项D错误。} 计算结果不正确，应为4.267。
\end{explanation}
\checklist{超阈值方法预期损失计算：ES = VaR/(1-ξ) + (β-ξu)/(1-ξ)}


\begin{question}
    Extreme value theory (EVT) can assist with value at risk (VaR) calculations by providing better probability estimates of observing extreme losses than that indicated by a standard normal distribution because empirical distributions exhibit fat tails. If one uses the generalized Pareto distribution (GPD) method to generate parameter estimates for the shape parameter, fat tails will indicate a:
\end{question}

\begin{choices}
    \item positive parameter estimate and VaR calculations that are too large
    \item negative parameter estimate and VaR calculations that are too small
    \item positive parameter estimate and VaR calculations that are too small
    \item negative parameter estimate and VaR calculations that are too large
\end{choices}

\begin{correctanswer}
    C
\end{correctanswer}

\begin{explanation}
    \textbf{选项C正确。}
    
    数据中的厚尾产生正形状参数(ξ > 0)。正形状参数表明分布比正态分布有更重的尾部。这意味着极端损失比正态分布预测的更可能发生。因此，基于正态分布的VaR估计偏小。
    
    \textbf{选项A错误。} 虽然厚尾确实表示正形状参数，但这实际上导致VaR估计偏小，而不是偏大。
    
    \textbf{选项B错误。} 厚尾表示正形状参数，而不是负形状参数。负形状参数表示薄尾。
    
    \textbf{选项D错误。} 厚尾表示正形状参数，而不是负形状参数。负形状参数表示薄尾。
\end{explanation}
\checklist{极值理论：正形状参数表示厚尾，VaR估计偏小}

\begin{question}
    The peaks-over-threshold approach generally requires:
\end{question}

\begin{choices}
    \item More estimated parameters than the GEV approach and shares one parameter with the GEV
    \item Fewer estimated parameters than the GEV approach and shares one parameter with the GEV
    \item More estimated parameters than the GEV approach and does not share any parameters with the GEV approach
    \item Fewer estimated parameters than the GEV approach and does not share any parameters with the GEV approach
\end{choices}

\begin{correctanswer}
    B
\end{correctanswer}

\begin{explanation}
    \textbf{选项B正确。}
    
    超阈值方法通常需要三个参数：
    - 阈值(u)
    - 尺度参数(β)
    - 形状参数(ξ)
    
    广义极值方法需要更多参数。两种方法都共享形状参数(ξ)。对于金融数据，ξ > 0表示厚尾。
    
    \textbf{选项A错误。} 超阈值方法实际上比广义极值方法需要更少的参数，而不是更多。
    
    \textbf{选项C错误。} POT不需要比GEV更多的参数，两种方法都共享形状参数ξ。
    
    \textbf{选项D错误。} 虽然POT确实比GEV需要更少的参数，但两种方法都共享形状参数ξ。
\end{explanation}
\checklist{超阈值方法：参数更少但共享形状参数}


\begin{question}
    A researcher using the POT approach observes the following parameter values: β = 1.2, ξ = 0.20, u = 2.5\%, and Nu/n = 5\%. The 5\% VaR in percentage terms is:
\end{question}

\begin{choices}
    \item 1.456
    \item 2.234
    \item 2.856
    \item 18.745
\end{choices}

\begin{correctanswer}
    B
\end{correctanswer}

\begin{explanation}
    \textbf{选项B正确。}
    
    使用超阈值方法的VaR公式：
    \[
    \text{VaR} = u + \frac{\beta}{\xi} \left\{ \left[ \frac{N}{n} (1 - \text{confidence level}) \right]^{-\xi} - 1 \right\}
    \]
    
    代入给定值：
    \[
    \text{VaR} = 2.5 + \frac{1.2}{0.20} \left\{ \left[ \frac{1}{0.05 \times (1 - 0.95)} \right]^{-0.20} - 1 \right\} = 2.234\%
    \]
    
    其中：u(阈值) = 2.5\%，β(尺度参数) = 1.2，ξ(形状参数) = 0.20，N/n = 5\%，置信水平 = 95\%
    
    \textbf{选项A错误。} 计算结果不正确，应为2.234\%。
    
    \textbf{选项C错误。} 计算结果不正确，应为2.234\%。
    
    \textbf{选项D错误。} 计算结果不正确，应为2.234\%。
\end{explanation}
\checklist{超阈值方法VaR计算：包含阈值、尺度参数、形状参数}


\begin{question}
    The peaks-over-threshold (POT) approach is used by a firm to apply extreme value theory (EVT) to the distribution of excess losses over a high threshold. The firm estimated the following parameter values: distribution scale parameter = 1.10, distribution shape parameter = 0.20, threshold = 1.5\%, and number of observations that exceed threshold / threshold = 6\%. Compute the 1\% VaR in percentage terms and the corresponding expected shortfall measure.
\end{question}

\begin{choices}
    \item VaR = 3.25\%, and ES = 4.89\%
    \item VaR = 3.12\%, and ES = 4.45\%
    \item VaR = 2.98\%, and ES = 4.21\%
    \item VaR = 2.85\%, and ES = 4.12\%
\end{choices}

\begin{correctanswer}
    A
\end{correctanswer}

\begin{explanation}
    \textbf{选项A正确。}
    
    使用超阈值方法公式：
    \[
    \text{VaR} = u + \frac{\beta}{\xi} \left\{ \left[ \frac{N}{n} (1 - \text{confidence level}) \right]^{-\xi} - 1 \right\}
    \]
    \[
    \text{ES} = \frac{\text{VaR}}{1-\xi} + \frac{\beta - \xi u}{1-\xi}
    \]
    
    代入给定值：
    \[
    \text{VaR} = 1.5 + \frac{1.10}{0.20} \left\{ \left[ \frac{1}{0.06} (1 - 0.99) \right]^{-0.20} - 1 \right\} = 3.25\%
    \]
    \[
    \text{ES} = \frac{3.25}{1-0.20} + \frac{1.10-0.20 \times 1.5}{1-0.20} = 4.89\%
    \]
    
    其中：u(阈值) = 1.5\%，β(尺度参数) = 1.10，ξ(形状参数) = 0.20，N/n = 6\%，置信水平 = 99\%

\end{explanation}
\checklist{超阈值方法VaR和ES计算：包含阈值、尺度参数、形状参数}





\begin{question}
    A senior quantitative risk analyst is explaining to a group of junior
    analysts the importance of extreme-value theory(EVT) and how it is
    applied in the bank's risk management function. The senior analyst 
    identifies the Frechet, Gumbel, and Weibull distributions as special
    variations of the generalized extreme-value distribution, and describes the 
    location, scale, and tail index parameters for each distribution. Which of 
    the following distributions is consistent with a tail index parameter (ξ) 
    whose value is greater than zero？
    
\end{question}

\begin{choices}
    \item Weibull
    \item Gumbel
    \item Fréchet
    \item Binomial
\end{choices}

\begin{correctanswer}
    C
\end{correctanswer}

\begin{explanation}
    \textbf{A选项错误。}
    \begin{itemize}
        \item Weibull分布对应ξ < 0的情况，其尾部比正态分布更轻
    \end{itemize}
    \textbf{B选项错误。}
    \begin{itemize}
        \item Gumbel分布对应ξ = 0的情况，其尾部与正态分布和对数正态分布类似
    \end{itemize}
    \textbf{C选项正确。}
    \begin{itemize}
        \item Fréchet分布对应ξ > 0的情况，具有重尾特征
        \item 这种分布适用于描述极端事件，如t分布和帕累托分布
        \item 在风险管理中，重尾分布更适合捕捉极端风险
    \end{itemize}
    \textbf{D选项错误。}
    \begin{itemize}
        \item 二项分布不属于极值理论中的广义极值分布族
    \end{itemize}
\end{explanation}
\checklist{
    掌握极值理论中三种分布的特征及其应用场景，理解尾部指数参数ξ与分布类型的关系
}

\begin{question}
    某投资银行的市场风险分析师正在研究如何使用极值理论(EVT)来建模低概率、高影响的事件。该分析师需要决定是使用广义极值理论(GEV)还是超阈值方法(POT)来建模损益(P\&L)分布的尾部。以下哪项陈述是正确的？
\end{question}

\begin{choices}
    \item POT方法提供了一个严格条件性的损失分布评估方法，因为极端损失通常受动态过程控制
    \item POT方法表明，增加多元分布中独立变量的数量可以生成大量极端观测值用于建模尾部分布
    \item GEV理论通过尾部指数参数来确定应该使用哪种分布来建模极端损失值，该参数反映了P\&L分布尾部的厚度
    \item GEV理论概述了将位置和尺度等中心趋势参数转换为可用于解释极端尾部损失分布的形式的程序
\end{choices}

\begin{correctanswer}
    C
\end{correctanswer}

\begin{explanation}
    \textbf{A选项错误。}
    \begin{itemize}
        \item POT方法不是严格条件性的，它关注的是超过阈值的观测值
    \end{itemize}
    \textbf{B选项错误。}
    \begin{itemize}
        \item POT方法并不依赖于增加独立变量数量，而是关注超过阈值的观测值
    \end{itemize}
    \textbf{C选项正确。}
    \begin{itemize}
        \item GEV理论确实通过尾部指数参数来确定合适的分布类型
        \item 这个参数反映了分布尾部的厚度，是选择分布类型的关键依据
        \item 这种方法在风险管理中具有重要的实践意义
    \end{itemize}
    \textbf{D选项错误。}
    \begin{itemize}
        \item GEV理论主要关注极值分布的类型选择，而不是参数转换
    \end{itemize}
\end{explanation}
\checklist{
    理解GEV理论的核心是通过尾部指数参数来确定合适的极值分布类型，这是建模极端风险的关键
}


\begin{question}
    某大型银行的高级风险经理被要求评估将广义极值理论(GEV)纳入银行风险计量流程的方法和影响。该经理正在研究计算极端损失分布VaR(EV-VaR)的程序，包括生成准确EV-VaR估计所需的假设和参数。以下哪项陈述是正确的？
\end{question}

\begin{choices}
    \item 错误地认为尾部指数参数在统计上不显著异于零，可能导致EV-VaR高估极端风险
    \item 由于选择正确的极值分布存在显著的模型风险，应选择最保守的Fréchet分布
    \item 尾部指数参数的值对任何类型GEV分布的分位数和EV-VaR估计都有显著影响
    \item 由于极值位于原始损失分布的尾部，在给定置信水平下生成的EV-VaR将提供极端风险的准确估计
\end{choices}

\begin{correctanswer}
    B
\end{correctanswer}

\begin{explanation}
    \textbf{A选项错误。}
    \begin{itemize}
        \item 错误地认为尾部指数参数不显著异于零会导致低估风险，而不是高估
    \end{itemize}
    \textbf{B选项正确。}
    \begin{itemize}
        \item 在存在模型风险的情况下，选择最保守的Fréchet分布是合理的
        \item Fréchet分布具有重尾特征，更适合捕捉极端风险
        \item 这种保守方法符合风险管理的基本原则
    \end{itemize}
    \textbf{C选项错误。}
    \begin{itemize}
        \item 尾部指数参数主要影响分布类型的选择，而不是直接影响VaR估计
    \end{itemize}
    \textbf{D选项错误。}
    \begin{itemize}
        \item 仅基于置信水平生成的EV-VaR不一定能准确反映极端风险
        \item 还需要考虑分布类型和参数估计的准确性
    \end{itemize}
\end{explanation}
\checklist{
    理解在模型风险存在的情况下，选择保守的Fréchet分布进行极端风险建模的重要性
}


\newpage




\section{考点5:VaR回测}
\setcounter{question}{0}

\begin{question}
    A portfolio manager is backtesting a sample at the 95\% confidence level to see if a VaR model needs to be recalibrated. He is using 300 daily returns for the sample and discovered 20 exceptions. What is the Z-Score for this sample when conducting VaR model verification?
\end{question}

\begin{choices}
    \item 0.85
    \item 1.45
    \item 1.89
    \item 3.12
\end{choices}

\begin{correctanswer}
    B
\end{correctanswer}

\begin{explanation}
    \textbf{选项B正确。}
    
    使用VaR回测的z值公式：
    \[
    Z = \frac{x - pN}{\sqrt{p(1-p)N}}
    \]
    
    代入给定值：
    \[
    Z = \frac{20 - 0.05 \times 300}{\sqrt{0.05 \times 0.95 \times 300}} = 1.45
    \]
    
    其中：x = 异常数量 = 20，p = 异常概率 = 0.05 (1 - 95\%)，N = 观测数量 = 300
    
    \textbf{选项A错误。} 计算结果不正确，应为1.45。
    
    \textbf{选项C错误。} 计算结果不正确，应为1.45。
    
    \textbf{选项D错误。} 计算结果不正确，应为1.45。
\end{explanation}
\checklist{VaR回测Z值计算：异常数量与期望值的标准化差异}


\begin{question}
    Basel II requires a backtest at a 99\% confidence level of a bank's internal value at risk (VaR) model (IMA). Assume the bank's ten-day 99\% VaR is \$1.5 million (minimum of 99\% is hardwired per Basel). The null hypothesis is: the VaR model is accurate. Out of 1,200 observations, 30 exceptions are observed (we saw the actual loss exceed the VaR 30 out of 1200 observations).
\end{question}

\begin{choices}
    \item We will probably call the VaR model good (accurate) but we risk a Type I error.
    \item We will probably call the VaR model good (accurate) but we risk a Type II error.
    \item We will probably call the model bad (inaccurate) but we risk a Type I error.
    \item We will probably call the model bad (inaccurate) but we risk a Type II error.
\end{choices}

\begin{correctanswer}
    C
\end{correctanswer}

\begin{explanation}
    \textbf{选项C正确。}
    
    期望异常数 = 1200 × (1-0.99) = 12
    实际异常数 = 30
    
    由于30 > 12，我们可能会拒绝模型。拒绝原假设(模型准确)时，当原假设可能为真时，这是I型错误。
    
    \textbf{选项A错误。} 在1200个观测值中有30个异常，模型可能会被拒绝。I型错误发生在拒绝真实原假设时。
    
    \textbf{选项B错误。} 在1200个观测值中有30个异常，模型可能会被拒绝。II型错误发生在未能拒绝错误原假设时。
    
    \textbf{选项D错误。} II型错误发生在未能拒绝错误原假设时。在这种情况下，我们正在拒绝原假设。
\end{explanation}
\checklist{VaR回测假设检验：拒绝原假设时存在I型错误风险}


\begin{question}
    Which of the following statements regarding verification of a VaR model by examining its failure rates is false?
\end{question}

\begin{choices}
    \item The frequency of exceptions should correspond to the confidence level used for the model.
    \item According to Kupiec (1995), we should reject the hypothesis that the model is correct if the LF>3.84.
    \item Backtesting VaR models with lower confidence levels is difficult because the number of exceptions is not high enough to provide meaningful information.
    \item The range for the number of exceptions must strike a balance between the chances of rejecting an accurate model (a Type I error) and the chance of accepting an inaccurate model (a Type II error).
\end{choices}

\begin{correctanswer}
    C
\end{correctanswer}

\begin{explanation}
    \textbf{选项C正确。}
    
    这个陈述是错误的，因为实际上是更高的置信水平使回测变得困难。在更高置信水平(如99\%)下，异常数量较少。较少的异常使得难以得出统计显著的结论。较低置信水平(如95\%)为测试提供更多异常数据。
    
    \textbf{选项A错误。} 异常频率应该与模型使用的置信水平匹配。例如，在95\%置信水平下，我们期望5\%的异常。
    
    \textbf{选项B错误。} Kupiec检验使用卡方分布。临界值3.84对应于95\%置信水平。
    
    \textbf{选项D错误。} 范围必须在I型错误和II型错误之间取得平衡。范围太窄会增加I型错误风险，范围太宽会增加II型错误风险。
\end{explanation}
\checklist{VaR回测：高置信度导致异常少，低置信度提供更多测试数据}


\begin{question}
    A portfolio manager is analyzing a 1-day 97\% VaR model. Assuming 252 days in a year, what is the maximum number of daily losses exceeding the 1-day 97\% VaR that is acceptable in a 1-year backtest to conclude, at a 95\% confidence level, that the model is calibrated correctly?
\end{question}

\begin{choices}
    \item 7
    \item 12
    \item 15
    \item 18
\end{choices}

\begin{correctanswer}
    B
\end{correctanswer}

\begin{explanation}
    \textbf{选项B正确。}
    
    使用VaR回测的z值公式：
    \[
    \frac{x - pT}{\sqrt{p(1-p)T}} \leq 1.96
    \]
    
    代入给定值：
    \[
    x \leq 1.96 \times \sqrt{0.03 \times 0.97 \times 252} + 0.03 \times 252 = 12.35
    \]
    
    其中：p = 失败率 = 1 - 97\% = 3\%，T = 观测数量 = 252，1.96 = 95\%置信水平临界值
    
    因此，最大可接受异常数 = 12
    
    \textbf{选项A错误。} 与计算的最大值12相比太低。
    
    \textbf{选项C错误。} 超过了计算的最大值12个异常。
    
    \textbf{选项D错误。} 超过了计算的最大值12个异常。
\end{explanation}
\checklist{VaR回测最大异常数：x ≤ 1.96×√[p(1-p)T] + pT}


\begin{question}
    The Basel Committee has established four categories of causes for exceptions. Which of the following does not apply to one of those categories?
\end{question}

\begin{choices}
    \item The sample is small.
    \item Intraday trading activity.
    \item Model accuracy needs improvement.
    \item The basic integrity of the model is lacking.
\end{choices}

\begin{correctanswer}
    A
\end{correctanswer}

\begin{explanation}
    \textbf{选项A正确。}
    
    样本大小不是巴塞尔委员会规定的四个类别之一。四个官方类别是：
    - 运气不佳(随机机会)
    - 日内交易活动
    - 模型准确性需要提高
    - 模型基本完整性不足
    
    \textbf{选项B错误。} 日内交易活动是四个官方类别之一。
    
    \textbf{选项C错误。} 模型准确性需要提高是四个官方类别之一。
    
    \textbf{选项D错误。} 模型基本完整性不足是四个官方类别之一。
\end{explanation}
\checklist{巴塞尔VaR异常四类原因：运气不佳、日内交易、准确性、完整性}


\begin{question}
    Alpha Capital Management decided to test the VaR model for its Fixed Income Portfolio over 12 trading days in October. Trading on the portfolio ceased during the twelve days of the test. The backtesting provided the following predicted daily percentage returns for the portfolio, at a 95\% confidence level:

    \begin{table}[H]
        \centering
        \begin{tabular}{cccc}
            \hline
            Upper Limit &       &       &       \\
            \hline
            0.45 & -0.25 & 0.32 \\
            0.48 & -0.18 & 0.41 \\
            0.49 & -0.12 & -0.15 \\
            0.49 & -0.20 & -0.28 \\
            0.49 & -0.25 & 0.32 \\
            0.48 & -0.19 &       \\
            \hline
        \end{tabular}
    \end{table}

    Which statement about the backtesting is TRUE?
\end{question}

\begin{choices}
    \item At this confidence level, 18\% of the actual results in the backtesting should fall outside the predicted range.
    \item The backtesting is not valid because the portfolio was not traded during the twelve days. It is only appropriate to backtest a portfolio that has undergone significant changes between calculations because otherwise the results are not relevant to actual trading situations.
    \item At the end of the period, Alpha Capital will recompute VaR to determine what the upper and lower ranges of the portfolio return should have been.
\item The backtesting does not confirm the validity of the model.
\end{choices}

\begin{correctanswer}
    D
\end{correctanswer}

\begin{explanation}
    \textbf{选项D正确。}
    
    在95\%置信水平下，期望异常数 = 12 × 5\% = 0.6
    最大可接受异常数 = 1
    实际异常数 = 2(超出预测范围)
    
    由于2 > 1，模型未能通过回测。
    
    \textbf{选项A错误。} 在95\%置信水平下，只有5\%的结果应该超出范围。18\%远高于期望的5\%。
    
    \textbf{选项B错误。} 回测实际上更适合于投资组合未发生重大变化的情况。重大变化会使回测结果不太可靠。
    
    \textbf{选项C错误。} VaR范围在回测期间之前就已确定。在期间结束后重新计算会使回测过程无效。
\end{explanation}
\checklist{VaR回测：95\%置信水平下异常数不应超过1个}

\begin{question}
    When backtesting a VaR model. How many exceptions are expected when the model using a 98\% confidence level over a 252-day period?
\end{question}

\begin{choices}
    \item 3.5
    \item 5.0
    \item 6.5
    \item 8.0
\end{choices}

\begin{correctanswer}
    B
\end{correctanswer}

\begin{explanation}
    \textbf{选项B正确。}
    
    预期异常数 = (1 - 置信水平) × 天数
    = (1 - 0.98) × 252
    = 0.02 × 252
    = 5.04
    
    \textbf{选项A错误。} 3.5个异常对应于99.6\%置信水平。在98\%置信水平下，我们期望更多异常。
    
    \textbf{选项C错误。} 6.5个异常对应于97.4\%置信水平。在98\%置信水平下，我们期望更少异常。
    
    \textbf{选项D错误。} 8.0个异常对应于96.8\%置信水平。在98\%置信水平下，我们期望更少异常。
\end{explanation}
\checklist{VaR回测预期异常数：(1-置信水平)×天数}

\begin{question}
    A hedge fund conducted a backtest of its 98\% daily value at risk (VaR) and observed 8 exceptions over the last year which included 252 trading days (T = 252). If we use the normal distribution to approximate the binomial for purposes of model verification, what is our accept/reject opinion of the model under a 90\% two-tailed test?
\end{question}

\begin{choices}
    \item Accept with a test statistic of 1.45
    \item Accept with a test statistic of 2.12
    \item Reject with a test statistic of 1.45
    \item Reject with a test statistic of 2.12
\end{choices}

\begin{correctanswer}
    D
\end{correctanswer}

\begin{explanation}
    \textbf{选项D正确。}
    
    期望异常数 = (1 - 98\%) × 252 = 5.04
    标准误差 = √[p(1-p)T] = √[2\% × 98\% × 252] = 2.22
    检验统计量 = (8 - 5.04) / 2.22 = 2.12
    90\%双尾检验临界值 = 1.645
    
    由于2.12 > 1.645，我们拒绝原假设。模型被认为存在问题。
    
    \textbf{选项A错误。} 计算的检验统计量是2.12，不是1.45。即使使用1.45，我们仍会拒绝模型。
    
    \textbf{选项B错误。} 虽然检验统计量正确(2.12)，但我们应该拒绝模型，而不是接受它。
    
    \textbf{选项C错误。} 计算的检验统计量是2.12，不是1.45。拒绝决定是正确的。
\end{explanation}
\checklist{VaR回测假设检验：检验统计量超过临界值1.645时拒绝原假设}


\begin{question}
    Which of the following statements regarding verification of a VaR model by examining its failure rates is false?
\end{question}

\begin{choices}
    \item The frequency of exceptions should correspond to the confidence level used for the model.
    \item According to Kupiec (1995), we should reject the hypothesis that the model is correct if the LF>3.84.
    \item Backtesting VaR models with lower confidence levels is difficult because the number of exceptions is not high enough to provide meaningful information.
    \item The range for the number of exceptions must strike a balance between the chances of rejecting an accurate model (a Type I error) and the chance of accepting an inaccurate model (a Type II error).
\end{choices}

\begin{correctanswer}
    C
\end{correctanswer}

\begin{explanation}
    \textbf{选项C正确。}
    
    这个陈述是错误的，因为实际上是更高的置信水平使回测变得困难。在更高置信水平(如99\%)下，异常数量较少。较少的异常使得难以得出统计显著的结论。较低置信水平(如95\%)为测试提供更多异常数据。
    
    \textbf{选项A错误。} 异常频率应该与模型使用的置信水平匹配。例如，在99\%置信水平下，我们期望1\%的异常。
    
    \textbf{选项B错误。} Kupiec检验使用卡方分布。临界值3.84对应于95\%置信水平。
    
    \textbf{选项D错误。} 范围必须在I型错误和II型错误之间取得平衡。范围太窄会增加I型错误风险，范围太宽会增加II型错误风险。
\end{explanation}
\checklist{VaR回测：高置信度导致异常少，低置信度提供更多测试数据}


\begin{question}
    A risk manager is backtesting a daily holding period VaR model using a 99\% confidence level over a 252-day period and is using a 3.84 test statistic. The following table shows the calculated values of a log-likelihood ratio (LR) at a 99\% confidence level.

    \begin{table}[H]
        \centering
        \begin{tabular}{|c|cccccccccccc|}
            \hline
            \multicolumn{13}{|c|}{Number of Exceptions} \\
            \hline
            & 1 & 2 & 3 & 4 & 5 & 6 & 7 & 8 & 9 & 10 & 11 & 12 \\
            \hline
            & 7.25 & 4.38 & 2.56 & 1.33 & 0.62 & 0.45 & 0.59 & 0.92 & 1.51 & 2.34 & 3.37 & 4.59 \\
            \hline
        \end{tabular}
    \end{table}

    Based on the above information, which of the following statements accurately describes the VaR model that is being backtested?
\end{question}

\begin{choices}
    \item If the number of exceptions is more than 2, we would not reject the model.
    \item If the number of exceptions is more than 1 and less than 11, we may commit a Type II error.
    \item If the number of exceptions is less than 1, we would accept the hypothesis that the model is correct.
    \item If the number of exceptions is less than 1, we may commit a Type II error.
\end{choices}

\begin{correctanswer}
    B
\end{correctanswer}

\begin{explanation}
    \textbf{选项B正确。}
    
    当1 < 异常数 < 11时，LR < 3.84，我们不会拒绝模型。这可能导致II型错误(接受错误的模型)。II型错误发生在未能拒绝错误原假设时。
    
    \textbf{选项A错误。} 当异常数 > 2时，LR > 3.84，我们会拒绝模型，而不是接受它。
    
    \textbf{选项C错误。} 当异常数 < 1时，LR > 3.84，我们会拒绝模型，而不是接受它。
    
    \textbf{选项D错误。} 当异常数 < 1时，LR > 3.84，我们会拒绝模型。II型错误发生在未能拒绝错误模型时。
\end{explanation}
\checklist{VaR回测假设检验：LR < 3.84时不拒绝模型但可能犯II型错误}


\begin{question}
    You are given the following information about the returns of stock X and stock Y: Variance of return of stock X = 169.0. Variance of return of stock Y = 289.0. Covariance between the return of stock X and the return of stock Y = 85.0. At the end of 2024, you are holding USD 6 million in stock X. You are considering a strategy of shifting USD 2.5 million into stock Y and keeping USD 3.5 million in stock X. What percentage of risk, as measured by standard deviation of return, can be reduced by this strategy?
\end{question}

\begin{choices}
    \item 0.012
    \item 0.075
    \item 0.098
    \item 0.125
\end{choices}

\begin{correctanswer}
    B
\end{correctanswer}

\begin{explanation}
    \textbf{选项B正确。}
    
    原投资组合方差 = 6² × 169 = 6,084
    原投资组合波动率 = √6,084 = 78.0
    
    新投资组合方差 = 3.5² × 169 + 2.5² × 289 + 2 × 85.0 × 3.5 × 2.5 = 5,640
    新投资组合波动率 = √5,640 = 75.1
    
    风险降低 = (78.0 - 75.1) / 78.0 = 7.5\%
    
\end{explanation}
\checklist{投资组合风险计算：组合方差 = w₁²σ₁² + w₂²σ₂² + 2w₁w₂Cov₁₂}


\begin{question}
    Based on Basel III rules for backtesting, a penalty is given to banks that have more than four exceptions to their 1-day 99\% VaR over the course of 252 trading days. The supervisor gives these penalties based on four criteria. Which of the following causes of exceptions is most likely to lead to a penalty?
\end{question}

\begin{choices}
    \item The bank increases its intraday trading activity.
    \item A large move in exchange rates was combined with a small move in correlations.
    \item The bank's model calculates currency risk based on the average duration of the forwards in the portfolio.
    \item A sudden market crisis in a developed market leads to losses in the bond positions in that country.
\end{choices}

\begin{correctanswer}
    C
\end{correctanswer}

\begin{explanation}
    \textbf{选项C正确。}
    
    使用平均期限而不是适当的风险度量是明显的模型缺陷。模型本可以通过更好的方法学得到改进。巴塞尔要求银行使用适当的风险计量技术。
    
    \textbf{选项A错误。} 增加的日内交易活动可能导致异常，但这不一定是模型缺陷。这更多与交易行为相关，而不是模型准确性。
    
    \textbf{选项B错误。} 汇率大幅波动与相关性小幅变化相结合代表"运气不佳"或市场条件，而不是可以预防的模型缺陷。
    
    \textbf{选项D错误。} 发达市场的突然市场危机代表"运气不佳"或市场条件，而不是可以预防的模型缺陷。
\end{explanation}
\checklist{巴塞尔回测惩罚：模型缺陷最可能导致惩罚}

\begin{question}
    Basel III requires a backtest of a bank's internal value at risk (VaR) model (IMA). Assume the bank's ten-day 99\% VaR is \$3.5 million (minimum of 99\% is hard-wired per Basel). The null hypothesis is: the VaR model is accurate. Out of 800 observations, 12 exceptions are observed (we saw the actual loss exceed the VaR 12 out of 800 observations). (Binomial CDF)
\end{question}

\begin{choices}
    \item We will probably call the VaR model good (accurate) but we risk a Type I error.
    \item We will probably call the VaR model good (accurate) but we risk a Type II error.
    \item We will probably call the model bad (inaccurate) but we risk a Type I error.
    \item We will probably call the model bad (inaccurate) but we risk a Type II error.
\end{choices}

\begin{correctanswer}
    C
\end{correctanswer}

\begin{explanation}
    \textbf{A选项错误。}
    \begin{itemize}
        \item 此选项错误，因为：
        \begin{itemize}
            \item 800次观察中出现12次例外，模型很可能被拒绝
            \item I型错误发生在拒绝真实原假设时
        \end{itemize}
    \end{itemize}
    \textbf{B选项错误。}
    \begin{itemize}
        \item 此选项错误，因为：
        \begin{itemize}
            \item 800次观察中出现12次例外，模型很可能被拒绝
            \item II型错误发生在未能拒绝错误原假设时
        \end{itemize}
    \end{itemize}
    \textbf{C选项正确。}
    \begin{itemize}
        \item 此选项正确，因为：
        \begin{itemize}
            \item 预期例外数 = 800 × (1-0.99) = 8
            \item 实际例外数 = 12
            \item 出现12次或更多例外的概率很低
            \item 我们拒绝原假设（模型准确）
            \item 拒绝真实原假设是I型错误
        \end{itemize}
    \end{itemize}
    \textbf{D选项错误。}
    \begin{itemize}
        \item 此选项错误，因为：
        \begin{itemize}
            \item II型错误发生在未能拒绝错误原假设时
            \item 在此情况下，我们正在拒绝原假设
        \end{itemize}
    \end{itemize}
\end{explanation}
\checklist{
    掌握VaR回测中的假设检验：当实际例外数显著超过预期时，应拒绝原假设，此时存在I型错误风险（错误地拒绝了真实的原假设）
}


\begin{question}
    Which of the following statements regarding verification of a VAR model by examining its failure rates is false?I. The frequency of exceptions should correspond to the confidence level used for the model.II. According to Kupiec (1995), we should reject the hypothesis that the model is correct if the LR>3.84.III. Backtesting VAR models with higher confidence levels is difficult because the number of exceptions is not high enough to provide meaningful information.IV. The range for the number of exceptions must strike a balance between the chances of rejecting an accurate model (a type 1 error) and the chance of accepting an inaccurate model (a type 2 error)
\end{question}

\begin{choices}
    \item I and IV
    \item II only
    \item III only
    \item II and IV
\end{choices}

\begin{correctanswer}
    C
\end{correctanswer}

\begin{explanation}
    \textbf{A选项错误。}
    \begin{itemize}
        \item 此选项错误，因为：
        \begin{itemize}
            \item I正确：例外频率应与模型使用的置信水平相对应
            \item IV正确：例外数范围必须在拒绝准确模型（I型错误）和接受不准确模型（II型错误）之间取得平衡
        \end{itemize}
    \end{itemize}
    \textbf{B选项错误。}
    \begin{itemize}
        \item 此选项错误，因为：
        \begin{itemize}
            \item II正确：Kupiec检验使用3.84作为临界值
            \item III错误：是高置信水平使回测变得困难
        \end{itemize}
    \end{itemize}
    \textbf{C选项正确。}
    \begin{itemize}
        \item 此选项正确，因为：
        \begin{itemize}
            \item III错误：是高置信水平使回测变得困难
            \item 在高置信水平下（如99\%），例外数量较少
            \item 较少的例外使得难以得出统计上显著的结论
            \item 低置信水平（如95\%）提供更多例外用于测试
        \end{itemize}
    \end{itemize}
    \textbf{D选项错误。}
    \begin{itemize}
        \item 此选项错误，因为：
        \begin{itemize}
            \item II正确：Kupiec检验使用3.84作为临界值
            \item IV正确：例外数范围必须在I型和II型错误之间取得平衡
        \end{itemize}
    \end{itemize}
\end{explanation}
\checklist{
    掌握VaR回测中的关键概念：高置信度导致例外数量少，难以进行有效的回测；而低置信度则提供更多例外数据用于测试
}

\begin{question}
    An analyst is backtesting a daily holding period VaR model using a 99\% confidence level over a 240-day period and is using a 3.84 test statistic. The following table shows the calculated values of a log-likelihood ratio (LR) at a 99\% confidence level. 

    \begin{table}[H]
        \centering
        \begin{tabular}{|c|c|}
            \hline
            Exceptions & LR Value \\
            \hline
            1 & 0.6123 \\
            2 & 1.2345 \\
            3 & 1.8765 \\
            4 & 2.5432 \\
            5 & 3.1234 \\
            \hline
        \end{tabular}
    \end{table}
    Based on the above information, which of the following statements accurately describes the VaR model that is being backtested?
\end{question}

\begin{choices}
    \item If the number of exceptions is more than 3, we would not reject the model.
    \item If the number of exceptions is more than 1 and less than 10, we may commit a Type II error.
    \item If the number of exceptions is less than 1, we would accept the hypothesis that the model is correct.
    \item If the number of exceptions is less than 1, we may commit a Type II error.
\end{choices}

\begin{correctanswer}
    B
\end{correctanswer}

\begin{explanation}
    \textbf{A选项错误。}
    \begin{itemize}
        \item 此选项错误，因为：
        \begin{itemize}
            \item 当例外数 > 3时，LR > 3.84
            \item 我们会拒绝模型，而不是接受它
        \end{itemize}
    \end{itemize}
    \textbf{B选项正确。}
    \begin{itemize}
        \item 此选项正确，因为：
        \begin{itemize}
            \item 当1 < 例外数 < 10时，LR < 3.84
            \item 我们不会拒绝模型
            \item 这可能导致II型错误（接受错误的模型）
            \item II型错误发生在未能拒绝错误原假设时
        \end{itemize}
    \end{itemize}
    \textbf{C选项错误。}
    \begin{itemize}
        \item 此选项错误，因为：
        \begin{itemize}
            \item 当例外数 < 1时，LR > 3.84
            \item 我们会拒绝模型，而不是接受它
        \end{itemize}
    \end{itemize}
    \textbf{D选项错误。}
    \begin{itemize}
        \item 此选项错误，因为：
        \begin{itemize}
            \item 当例外数 < 1时，LR > 3.84
            \item 我们会拒绝模型
            \item II型错误发生在未能拒绝错误模型时
        \end{itemize}
    \end{itemize}
\end{explanation}
\checklist{
    掌握VaR回测中的假设检验：当LR < 3.84时不拒绝模型，但可能犯II型错误（接受错误的模型）
}


\begin{question}
    A risk analyst at Bank XYZ has been asked to verify the bank's VaR model based on the failure rate. Upon examining the bank's daily trading profits and losses over the past year, the analyst finds that there were 18 exceptions to the bank's 99\% VaR model. To determine whether the VaR model is biased, the analyst calculates the z-statistic for the model. Assuming there are 250 trading days in a year, which of the following is the correct estimate of the z-statistic?
\end{question}

\begin{choices}
    \item 2.15
    \item 2.23
    \item 2.31
    \item 2.42
\end{choices}

\begin{correctanswer}
    D
\end{correctanswer}

\begin{explanation}
    \textbf{A选项错误。}
    \begin{itemize}
        \item 此选项错误，因为：
        \begin{itemize}
            \item 计算得出的z值为2.42，不是2.15
            \item 计算错误可能出现在标准误差项
        \end{itemize}
    \end{itemize}
    \textbf{B选项错误。}
    \begin{itemize}
        \item 此选项错误，因为：
        \begin{itemize}
            \item 计算得出的z值为2.42，不是2.23
            \item 计算错误可能出现在标准误差项
        \end{itemize}
    \end{itemize}
    \textbf{C选项错误。}
    \begin{itemize}
        \item 此选项错误，因为：
        \begin{itemize}
            \item 计算得出的z值为2.42，不是2.31
            \item 计算错误可能出现在标准误差项
        \end{itemize}
    \end{itemize}
    \textbf{D选项正确。}
    \begin{itemize}
        \item 此选项正确，因为：
        \begin{align*}
            z &= \frac{18 - 250 \times 0.01}{\sqrt{250 \times 0.01 \times (1 - 0.01)}} \\
            &= \frac{18 - 2.5}{\sqrt{2.5 \times 0.99}} \\
            &= \frac{15.5}{\sqrt{2.475}} \\
            &= \frac{15.5}{1.573} \\
            &= 2.42
        \end{align*}
        \item 其中：
        \begin{itemize}
            \item 预期例外数 = 250 × 1\% = 2.5
            \item 标准误差 = √[250 × 1\% × 99\%] = 1.573
            \item z值 = (18 - 2.5) / 1.573 = 2.42
        \end{itemize}
    \end{itemize}
\end{explanation}
\checklist{
    掌握VaR回测中z值的计算：z = (x-pT)/√[p(1-p)T]，其中x为例外次数，p为例外概率，T为观察次数
}


\begin{question}
    A risk analyst backtests an internal 1-day 99\% VaR model over a 240-day look-back period and finds the model's failure rate to be 2.8\%. If the analyst conducts a one-tailed test at the 99\% confidence level and the number of exceptions follows a binomial probability distribution, what should the analyst conclude?
\end{question}

\begin{choices}
    \item The model is incorrectly calibrated since the test statistic is less than the critical value.
    \item The model is incorrectly calibrated since the test statistic is greater than the critical value.
    \item The model is correctly calibrated since the test statistic is less than the critical value.
    \item The model is correctly calibrated since the test statistic is greater than the critical value.
\end{choices}

\begin{correctanswer}
    B
\end{correctanswer}

\begin{explanation}
    \textbf{A选项错误。}
    \begin{itemize}
        \item 此选项错误，因为：
        \begin{itemize}
            \item 检验统计量大于临界值
            \item 模型校准不正确
        \end{itemize}
    \end{itemize}
    \textbf{B选项正确。}
    \begin{itemize}
        \item 此选项正确，因为：
        \begin{itemize}
            \item 预期失败率 = 1\%
            \item 实际失败率 = 2.8\%
            \item 例外数 = 240 × 2.8\% = 6.72
            \item 预期例外数 = 240 × 1\% = 2.4
            \item 检验统计量 = (6.72 - 2.4) / √[240 × 1\% × 99\%] = 2.78
            \item 99\%置信度下的临界值 = 2.326
            \item 由于2.78 > 2.326，我们拒绝原假设
            \item 模型校准不正确
        \end{itemize}
    \end{itemize}
    \textbf{C选项错误。}
    \begin{itemize}
        \item 此选项错误，因为：
        \begin{itemize}
            \item 检验统计量大于临界值
            \item 模型校准不正确
        \end{itemize}
    \end{itemize}
    \textbf{D选项错误。}
    \begin{itemize}
        \item 此选项错误，因为：
        \begin{itemize}
            \item 检验统计量大于临界值
            \item 模型校准不正确
        \end{itemize}
    \end{itemize}
\end{explanation}
\checklist{
    掌握VaR回测中的假设检验：当检验统计量超过临界值时，应拒绝原假设，认为模型校准不正确
}


\begin{question}
    A risk analyst backtests an internal 1-day 99\% VaR model over a 252-day look-back period and finds the model's failure rate to be 2.4\%. If the analyst conducts a one-tailed test at the 99\% confidence level and the number of exceptions follows a binomial probability distribution, what should the analyst conclude?
\end{question}

\begin{choices}
    \item The model is incorrectly calibrated since the test statistic is less than the critical value.
    \item The model is incorrectly calibrated since the test statistic is greater than the critical value.
    \item The model is correctly calibrated since the test statistic is less than the critical value.
    \item The model is correctly calibrated since the test statistic is greater than the critical value.
\end{choices}

\begin{correctanswer}
    B
\end{correctanswer}

\begin{explanation}
    \textbf{A选项错误。}
    \begin{itemize}
        \item This option is incorrect because:
    \begin{itemize}
            \item The test statistic is greater than the critical value
            \item The model is incorrectly calibrated
    \end{itemize}
    \end{itemize}
    \textbf{B选项正确。}
    \begin{itemize}
        \item This option is correct because:
    \begin{itemize}
            \item Expected failure rate = 1\%
            \item Actual failure rate = 2.4\%
            \item Number of exceptions = 252 × 2.4\% = 6.05
            \item Expected exceptions = 252 × 1\% = 2.52
            \item Test statistic = (6.05 - 2.52) / √[252 × 1\% × 99\%] = 2.23
            \item Critical value at 99\% confidence = 2.326
            \item Since 2.23 > 2.326, we reject the null hypothesis
            \item The model is incorrectly calibrated
    \end{itemize}
    \end{itemize}
    \textbf{C选项错误。}
    \begin{itemize}
        \item This option is incorrect because:
    \begin{itemize}
            \item The test statistic is greater than the critical value
            \item The model is incorrectly calibrated
    \end{itemize}
    \end{itemize}
    \textbf{D选项错误。}
    \begin{itemize}
        \item This option is incorrect because:
    \begin{itemize}
            \item The test statistic is greater than the critical value
            \item The model is incorrectly calibrated
    \end{itemize}
    \end{itemize}
\end{explanation}
\checklist{
    掌握VaR回测中的假设检验：当检验统计量超过临界值时，应拒绝原假设，认为模型校准不正确
}


\newpage




\section{考点6:基于超额收益的VaR回测}
\setcounter{question}{0}
\begin{question}
    Unconditional testing does not reflect:
\end{question}

\begin{choices}
    \item the size of the portfolio.
    \item the number of exceptions.
    \item the confidence level chosen.
    \item the timing of the exceptions.
\end{choices}

\begin{correctanswer}
    D
\end{correctanswer}

\begin{explanation}
    \textbf{A选项错误。}
    \begin{itemize}
        \item 此选项错误，因为：
        \begin{itemize}
            \item 非条件检验确实反映投资组合规模
            \item 投资组合规模影响VaR计算
        \end{itemize}
    \end{itemize}
    \textbf{B选项错误。}
    \begin{itemize}
        \item 此选项错误，因为：
        \begin{itemize}
            \item 非条件检验确实反映例外数量
            \item 例外数量是检验的关键组成部分
        \end{itemize}
    \end{itemize}
    \textbf{C选项错误。}
    \begin{itemize}
        \item 此选项错误，因为：
        \begin{itemize}
            \item 非条件检验确实反映置信水平
            \item 置信水平决定预期例外数量
        \end{itemize}
    \end{itemize}
    \textbf{D选项正确。}
    \begin{itemize}
        \item 此选项正确，因为：
        \begin{itemize}
            \item 非条件检验只考虑例外总数
            \item 不考虑例外何时发生
            \item 无法检测例外的聚集性
            \item 无法识别例外是否独立
        \end{itemize}
    \end{itemize}
\end{explanation}
\checklist{
    掌握非条件检验的特点：只关注例外总数，不考虑例外发生的时间点，无法检测例外的聚集性和独立性
}


\begin{question}
    What is a key weakness of the value at risk (VaR) measure? VaR:
\end{question}

\begin{choices}
    \item does not consider the severity of losses in the tail of the returns distribution.
    \item is quite difficult to compute.
    \item is sub-additive.
    \item has an insufficient amount of back-testing data.
\end{choices}

\begin{correctanswer}
    A
\end{correctanswer}

\begin{explanation}
    \textbf{A选项正确。}
    \begin{itemize}
        \item This option is correct because:
        \begin{itemize}
            \item VaR only tells us the maximum loss at a given confidence level
            \item It does not provide information about losses beyond the VaR threshold
            \item It ignores the severity of extreme losses in the tail
            \item This is why Expected Shortfall (ES) is often used as a complement
        \end{itemize}
    \end{itemize}
    \textbf{B选项错误。}
    \begin{itemize}
        \item This option is incorrect because:
        \begin{itemize}
            \item VaR is relatively easy to compute
            \item Many software packages can calculate VaR
            \item The computation methods are well-established
        \end{itemize}
    \end{itemize}
    \textbf{C选项错误。}
    \begin{itemize}
        \item This option is incorrect because:
        \begin{itemize}
            \item VaR is actually super-additive, not sub-additive
            \item The sum of individual VaRs is often greater than portfolio VaR
        \end{itemize}
    \end{itemize}
    \textbf{D选项错误。}
    \begin{itemize}
        \item This option is incorrect because:
        \begin{itemize}
            \item Back-testing data availability is not a weakness of VaR itself
            \item It's a practical issue, not a theoretical limitation
        \end{itemize}
    \end{itemize}
\end{explanation}
\checklist{
    掌握VaR的主要缺点：不考虑尾部损失的严重程度，这是为什么需要预期短缺(ES)作为补充的原因
}

\begin{question}
    XYZ's pension plan has \$450 million in assets with an expected return of 12 percent. The last thirty monthly returns are given below. What is the 10 percent monthly probability VaR for XYZ's pension plan? 18.92\% -18.75\% 28.34\% 6.45\% -1.23\% 15.67\% 4.56\% 8.45\% 1.89\% 32.12\% 27.89\% 16.23\% 7.34\% -5.67\% -1.23\% 10.23\% 9.45\% 6.78\% 2.34\% 12.67\% 12.45\% 7.89\% 9.12\% 19.67\% 7.89\% 17.23\% 6.78\% 35.67\% 32.89\% 8.67\%
\end{question}

\begin{choices}
    \item \$2,250,000
    \item \$5,400,000
    \item \$3,105,000
    \item \$65,025,000
\end{choices}

\begin{correctanswer}
    C
\end{correctanswer}

\begin{explanation}
    \textbf{A选项错误。}
    \begin{itemize}
        \item 此选项错误，因为：
        \begin{itemize}
            \item 计算得出的VaR是\$3,105,000，不是\$2,250,000
            \item 计算错误可能出现在收益率排序上
        \end{itemize}
    \end{itemize}
    \textbf{B选项错误。}
    \begin{itemize}
        \item 此选项错误，因为：
        \begin{itemize}
            \item 计算得出的VaR是\$3,105,000，不是\$5,400,000
            \item 计算错误可能出现在收益率排序上
        \end{itemize}
    \end{itemize}
    \textbf{C选项正确。}
    \begin{itemize}
        \item 此选项正确，因为：
        \begin{itemize}
            \item 排序后的月收益率（从低到高）：
            \begin{itemize}
                \item -18.75\%, -5.67\%, -1.23\%, -1.23\%, 1.89\%, 2.34\%, 4.56\%, 6.78\%, 6.78\%, 6.45\%, 7.89\%, 7.89\%, 7.34\%, 8.45\%, 8.67\%, 9.12\%, 9.45\%, 10.23\%, 12.45\%, 12.67\%, 15.67\%, 16.23\%, 17.23\%, 18.92\%, 19.67\%, 27.89\%, 28.34\%, 32.12\%, 32.89\%, 35.67\%
            \end{itemize}
            \item 10\%最低收益率是第3个值，即-1.23\%
            \item 因此，投资组合的10\% VaR = 0.0123 × \$450,000,000 = \$3,105,000
        \end{itemize}
    \end{itemize}
    \textbf{D选项错误。}
    \begin{itemize}
        \item 此选项错误，因为：
        \begin{itemize}
            \item 计算得出的VaR是\$3,105,000，不是\$65,025,000
            \item 计算错误可能出现在使用了错误的收益率
        \end{itemize}
    \end{itemize}
\end{explanation}
\checklist{
    掌握历史模拟法VaR的计算：将收益率从小到大排序，找到对应分位数的收益率，乘以资产价值得到VaR
}


\begin{question}
    In early 2024, a risk manager calculates the VaR for a hedge fund based on the last three years of data. The strategy of the fund is to buy stocks and write out-of-the-money calls. The manager needs to compute VaR. Which of the following methods would yield results that are least representative of the risks inherent in the portfolio?
\end{question}

\begin{choices}
    \item Historical simulation with full repricing
    \item Delta–normal VaR assuming zero drift
    \item Monte Carlo style VaR assuming zero drift with full repricing
    \item Historical simulation using delta-equivalents for all positions
\end{choices}

\begin{correctanswer}
    D
\end{correctanswer}

\begin{explanation}
    \textbf{A选项错误。}
    \begin{itemize}
        \item 此选项错误，因为：
        \begin{itemize}
            \item 完全重定价的历史模拟法适用于期权
            \item 它捕捉期权的非线性收益
            \item 使用实际历史价格变动
        \end{itemize}
    \end{itemize}
    \textbf{B选项错误。}
    \begin{itemize}
        \item 此选项错误，因为：
        \begin{itemize}
            \item Delta-正态VaR虽然简单，但可以作为初步近似
            \item 零漂移假设对风险测量是保守的
        \end{itemize}
    \end{itemize}
    \textbf{C选项错误。}
    \begin{itemize}
        \item 此选项错误，因为：
        \begin{itemize}
            \item 完全重定价的蒙特卡洛法适用于期权
            \item 它捕捉期权的非线性收益
            \item 零漂移假设是保守的
        \end{itemize}
    \end{itemize}
    \textbf{D选项正确。}
    \begin{itemize}
        \item 此选项正确，因为：
        \begin{itemize}
            \item 对期权使用delta等价法是不合适的
            \item 它忽略期权收益的非线性特性
            \item 低估虚值看涨期权的风险
            \item 投资组合包含需要完全重定价的期权
            \item Delta等价法对这个复杂投资组合过于简化
        \end{itemize}
    \end{itemize}
\end{explanation}
\checklist{
    掌握VaR计算方法的选择：对于包含期权的投资组合，应使用完全重定价方法，而不是简单的delta等价法
}

\begin{question}
    Consider a trader with an investment in a corporate bond with face value of \$300,000 and default probability of 0.6\%. Over the next period, we can either have no default, with a return of zero, or default with a loss of \$300,000. The payoffs are thus \$300,000 with probability of 0.6\% and +\$0 with probability of 99.4\%. Since the probability of getting \$0 is greater than 99\%, the VaR at the 99\% confidence level is \$0, without taking the mean into account. This is consistent with the definition that VaR is the smallest loss, such that the right-tail probability is at least 99\%. Now, consider a portfolio invested in four bonds (A, B, C, D) with the same characteristics and independent payoffs. Please compute the portfolio VaR at the 99\% confidence level (using loss distribution method):
\end{question}

\begin{choices}
    \item \$0
    \item \$300,000
    \item \$600,000
    \item \$900,000
\end{choices}

\begin{correctanswer}
    B
\end{correctanswer}

\begin{explanation}
    \textbf{A选项错误。}
    \begin{itemize}
        \item 此选项错误，因为：
        \begin{itemize}
            \item 对于四个债券，至少一个违约的概率更高
            \item 99\% VaR必须考虑一个违约的可能性
        \end{itemize}
    \end{itemize}
    \textbf{B选项正确。}
    \begin{itemize}
        \item 此选项正确，因为：
        \begin{itemize}
            \item 无违约概率 = 0.994⁴ = 0.976
            \item 一个违约概率 = 4 × 0.006 × 0.994³ = 0.0238
            \item 两个违约概率 = 6 × 0.006² × 0.994² = 0.000214
            \item 三个违约概率 = 4 × 0.006³ × 0.994 = 0.00000086
            \item 四个违约概率 = 0.006⁴ = 0.0000000013
            \item 在99\%置信水平下，我们需要考虑第一个违约
            \item 因此，VaR = \$300,000
        \end{itemize}
    \end{itemize}
    \textbf{C选项错误。}
    \begin{itemize}
        \item 此选项错误，因为：
        \begin{itemize}
            \item 两个违约的概率很低（0.0214\%）
            \item 超出了99\%置信水平
        \end{itemize}
    \end{itemize}
    \textbf{D选项错误。}
    \begin{itemize}
        \item 此选项错误，因为：
        \begin{itemize}
            \item 三个违约的概率极低（0.000086\%）
            \item 远超出99\%置信水平
        \end{itemize}
    \end{itemize}
\end{explanation}
\checklist{
    掌握组合VaR的计算：对于独立债券组合，需要考虑违约概率的分布，在给定置信水平下确定VaR
}

\newpage



\section{考点7:VaR映射}
\setcounter{question}{0}


\begin{question}
    Which of these statements regarding risk factor mapping approaches is/are correct? 
    
    I. Under the cash flow mapping approach, only the risk associated with the average maturity of a fixed income portfolio is mapped. 
    
    II. Cash flow mapping is the least precise method of risk mapping for a fixed-income portfolio. 
    
    III. Under the duration mapping approach, the risk of a bond is mapped to a zero-coupon bond of the same duration. 
    
    IV. Using more risk factors generally leads to better risk measurement but also requires more time to be devoted to the modeling process and risk computation.
\end{question}

\begin{choices}
    \item I and II
    \item I, III, and IV
    \item III and IV
    \item IV only
\end{choices}

\begin{correctanswer}
    C
\end{correctanswer}

\begin{explanation}
    \textbf{A选项错误。}
    \begin{itemize}
        \item 此选项错误，因为：
        \begin{itemize}
            \item I错误：现金流映射分别考虑每个支付
            \item II错误：现金流映射是最精确的方法
        \end{itemize}
    \end{itemize}
    \textbf{B选项错误。}
    \begin{itemize}
        \item 此选项错误，因为：
        \begin{itemize}
            \item I错误：现金流映射分别考虑每个支付
            \item III正确：久期映射使用零息债券
            \item IV正确：更多风险因子提高准确性但增加复杂性
        \end{itemize}
    \end{itemize}
    \textbf{C选项正确。}
    \begin{itemize}
        \item 此选项正确，因为：
        \begin{itemize}
            \item III正确：
            \begin{itemize}
                \item 久期映射使用零息债券
                \item 零息债券具有与原债券相同的久期
                \item 这在保持准确性的同时简化了风险测量
            \end{itemize}
            \item IV正确：
            \begin{itemize}
                \item 更多风险因子捕捉更多市场变动
                \item 这提高了风险测量的准确性
                \item 但需要更多计算资源和时间
            \end{itemize}
        \end{itemize}
    \end{itemize}
    \textbf{D选项错误。}
    \begin{itemize}
        \item 此选项错误，因为：
        \begin{itemize}
            \item III也正确：久期映射使用零息债券
            \item IV正确：更多风险因子提高准确性但增加复杂性
        \end{itemize}
    \end{itemize}
\end{explanation}
\checklist{
    掌握风险因子映射方法：现金流映射最精确，久期映射使用零息债券，风险因子越多测量越精确但计算越复杂
}

\begin{question}
    There is a short position in 1-year bonds with a \$250 million face value and a 6\% annual interest rate, with interest paid semiannually. The annualized interest rate on zero-coupon bonds is 4.2\% for a 6-month maturity and 4.5\% for a 12-month maturity. Decompose the bond into the cash flows of the two standard instruments, and then determine the present value of the cash flows of the standard instruments. What are the present values of each cash flow?
\end{question}

\begin{choices}
    \item (PV of CF1) -\$6,123,456; (PV of CF2) -\$233,456,789
    \item (PV of CF1) -\$6,234,567; (PV of CF2) -\$235,678,901
    \item (PV of CF1) -\$6,234,567; (PV of CF2) -\$240,384,615
    \item (PV of CF1) -\$6,456,789; (PV of CF2) -\$238,095,238
\end{choices}

\begin{correctanswer}
    C
\end{correctanswer}

\begin{explanation}
    \textbf{A选项错误。}
    \begin{itemize}
        \item 此选项错误，因为：
        \begin{itemize}
            \item 计算的现值不正确
            \item 计算错误可能出现在折现过程中
        \end{itemize}
    \end{itemize}
    \textbf{B选项错误。}
    \begin{itemize}
        \item 此选项错误，因为：
        \begin{itemize}
            \item 计算的现值不正确
            \item 计算错误可能出现在折现过程中
        \end{itemize}
    \end{itemize}
    \textbf{C选项正确。}
    \begin{itemize}
        \item 此选项正确，因为：
        \begin{itemize}
            \item 第一个现金流（6个月）：
            \begin{itemize}
                \item 息票支付 = \$250,000,000 × (0.06/2) = \$7,500,000
                \item 现值 = -\$7,500,000 / (1 + 0.042/2) = -\$6,234,567
            \end{itemize}
            \item 第二个现金流（12个月）：
            \begin{itemize}
                \item 息票 + 本金 = \$7,500,000 + \$250,000,000 = \$257,500,000
                \item 现值 = -\$257,500,000 / (1 + 0.045) = -\$240,384,615
            \end{itemize}
        \end{itemize}
    \end{itemize}
    \textbf{D选项错误。}
    \begin{itemize}
        \item 此选项错误，因为：
        \begin{itemize}
            \item 计算的现值不正确
            \item 计算错误可能出现在折现过程中
        \end{itemize}
    \end{itemize}
\end{explanation}
\checklist{
    掌握债券现金流分解和现值计算：将债券分解为标准工具（零息债券），使用相应的零息利率进行折现
}


\begin{question}
    Which of the following methods is not one of the three approaches for mapping a portfolio of fixed-income securities onto risk factors?
\end{question}

\begin{choices}
    \item Principal mapping
    \item Duration mapping
    \item Cash flow mapping
    \item Present value mapping
\end{choices}

\begin{correctanswer}
    D
\end{correctanswer}

\begin{explanation}
    \textbf{A选项错误。}
    \begin{itemize}
        \item 此选项错误，因为：
        \begin{itemize}
            \item 本金映射是三种标准方法之一
            \item 它将债券映射到具有相同期限的零息债券
            \item 这是最简单但最不准确的方法
        \end{itemize}
    \end{itemize}
    \textbf{B选项错误。}
    \begin{itemize}
        \item 此选项错误，因为：
        \begin{itemize}
            \item 久期映射是三种标准方法之一
            \item 它将债券映射到具有相同久期的零息债券
            \item 比本金映射更准确
        \end{itemize}
    \end{itemize}
    \textbf{C选项错误。}
    \begin{itemize}
        \item 此选项错误，因为：
        \begin{itemize}
            \item 现金流映射是三种标准方法之一
            \item 它将每个现金流映射到零息债券
            \item 是最准确但最复杂的方法
        \end{itemize}
    \end{itemize}
    \textbf{D选项正确。}
    \begin{itemize}
        \item 此选项正确，因为：
        \begin{itemize}
            \item 现值映射不是标准方法
            \item 三种标准方法是：
            \begin{itemize}
                \item 本金映射
                \item 久期映射
                \item 现金流映射
            \end{itemize}
        \end{itemize}
    \end{itemize}
\end{explanation}
\checklist{
    掌握固定收益证券风险因子映射的三种标准方法：本金映射、久期映射和现金流映射
}

\begin{question}
    A portfolio manager owns a portfolio of options on a non-dividend paying stock XYZ. The portfolio is made up of 20,000 deep in-the-money call options on XYZ and 50,000 deep out-of-the-money call options on XYZ. The portfolio also contains 30,000 forward contracts on XYZ. XYZ is trading at USD 150. If the volatility of XYZ is 30\% per year, which of the following amounts would be closest to the 1-day VaR of the portfolio at the 95 percent confidence level, assuming 250 trading days in a year?
\end{question}

\begin{choices}
    \item USD 2,340
    \item USD 234,000
    \item USD 280,800
    \item USD 318,240
\end{choices}

\begin{correctanswer}
    B
\end{correctanswer}

\begin{explanation}
    \textbf{A选项错误。}
    \begin{itemize}
        \item 此选项错误，因为：
        \begin{itemize}
            \item 计算的VaR太小
            \item 计算错误可能出现在delta近似上
        \end{itemize}
    \end{itemize}
    \textbf{B选项正确。}
    \begin{itemize}
        \item 此选项正确，因为：
        \begin{itemize}
            \item 步骤1：计算投资组合delta
            \begin{itemize}
                \item 深度实值看涨期权：delta ≈ 1
                \item 深度虚值看涨期权：delta ≈ 0
                \item 远期合约：delta = 1
                \item 总delta = 20,000 × 1 + 50,000 × 0 + 30,000 × 1 = 50,000
            \end{itemize}
            \item 步骤2：计算1天VaR
            \begin{itemize}
                \item VaR = α × S × Δ × σ × √(1/T)
                \item = 1.645 × 150 × 50,000 × 0.30 × √(1/250)
                \item = 234,000
            \end{itemize}
        \end{itemize}
    \end{itemize}
    \textbf{C选项错误。}
    \begin{itemize}
        \item 此选项错误，因为：
        \begin{itemize}
            \item 计算的VaR太高
            \item 计算错误可能出现在delta近似上
        \end{itemize}
    \end{itemize}
    \textbf{D选项错误。}
    \begin{itemize}
        \item 此选项错误，因为：
        \begin{itemize}
            \item 计算的VaR太高
            \item 计算错误可能出现在delta近似上
        \end{itemize}
    \end{itemize}
\end{explanation}
\checklist{
    掌握期权组合VaR的计算：使用delta-normal方法，考虑深度实值期权delta≈1，深度虚值期权delta≈0，远期delta=1
}


\begin{question}
    Computing VaR on a portfolio containing a very large number of positions can be simplified by mapping these positions to a smaller number of elementary risk factors. Which of the following mappings would be adequate?
\end{question}

\begin{choices}
    \item EUR/JPY forward contracts are mapped on the EUR/USD spot exchange rate.
    \item Each position in a corporate bond portfolio is mapped on the bond with the closest maturity among a set of government bonds.
    \item Government bonds paying regular coupons are mapped on zero-coupon government bonds.
    \item A position in the NASDAQ index is mapped on a position in a single stock within that index.
\end{choices}

\begin{correctanswer}
    C
\end{correctanswer}

\begin{explanation}
    \textbf{A选项错误。}
    \begin{itemize}
        \item 此选项错误，因为：
        \begin{itemize}
            \item EUR/JPY和EUR/USD是不同的货币对
            \item 它们有不同的风险因子和相关性
            \item 这种映射无法捕捉JPY敞口的特定风险
        \end{itemize}
    \end{itemize}
    \textbf{B选项错误。}
    \begin{itemize}
        \item 此选项错误，因为：
        \begin{itemize}
            \item 公司债券和政府债券有不同的信用风险
            \item 仅基于期限的映射忽略了信用利差风险
            \item 这会低估投资组合的风险
        \end{itemize}
    \end{itemize}
    \textbf{C选项正确。}
    \begin{itemize}
        \item 此选项正确，因为：
        \begin{itemize}
            \item 两种债券都是政府发行的
            \item 它们对利率风险有相似的敏感性
            \item 零息债券是风险测量的标准工具
            \item 映射虽然不完美，但对风险管理是适当的
        \end{itemize}
    \end{itemize}
    \textbf{D选项错误。}
    \begin{itemize}
        \item 此选项错误，因为：
        \begin{itemize}
            \item 单只股票无法代表整个指数
            \item 忽略了分散化收益
            \item 根据选择的股票会高估或低估风险
        \end{itemize}
    \end{itemize}
\end{explanation}
\checklist{
    掌握风险因子映射的原则：映射应该保持风险特征相似性，使用标准工具，并考虑风险因素的敏感性
}

\begin{question}
    An analyst is using the delta-normal method to determine the VAR of a fixed income portfolio. The portfolio contains a long position in 1-year bonds with a \$3 million face value and a 6\% coupon that is paid semiannually. The interest rates on 6- and 12-month maturity zero-coupon bonds are 4\% and 4.5\%, respectively. Mapping the long position to standard positions in the 6- and 12-month zeros, respectively, provides which of the following mapped positions?
\end{question}

\begin{choices}
    \item \$87,379 and \$2,949,276
    \item \$90,000 and \$3,090,000
    \item \$88,000 and \$2,970,000
    \item \$92,000 and \$3,120,000
\end{choices}

\begin{correctanswer}
    A
\end{correctanswer}

\begin{explanation}
    \textbf{A选项正确。}
    \begin{itemize}
        \item 此选项正确，因为：
        \begin{itemize}
            \item 步骤1：计算现金流
            \begin{itemize}
                \item 6个月息票 = \$3,000,000 × (0.06/2) = \$90,000
                \item 12个月支付 = \$90,000 + \$3,000,000 = \$3,090,000
            \end{itemize}
            \item 步骤2：计算现值
            \begin{itemize}
                \item 6个月零息债券：\$90,000 / (1 + 0.04/2) = \$87,379
                \item 12个月零息债券：\$3,090,000 / (1 + 0.045) = \$2,949,276
            \end{itemize}
        \end{itemize}
    \end{itemize}
    \textbf{B选项错误。}
    \begin{itemize}
        \item 此选项错误，因为：
        \begin{itemize}
            \item 显示的是未折现的现金流
            \item 计算错误出现在未应用现值
        \end{itemize}
    \end{itemize}
    \textbf{C选项错误。}
    \begin{itemize}
        \item 此选项错误，因为：
        \begin{itemize}
            \item 计算的现值不正确
            \item 计算错误可能出现在折现过程中
        \end{itemize}
    \end{itemize}
    \textbf{D选项错误。}
    \begin{itemize}
        \item 此选项错误，因为：
        \begin{itemize}
            \item 计算的现值不正确
            \item 计算错误可能出现在折现过程中
        \end{itemize}
    \end{itemize}
\end{explanation}
\checklist{
    掌握债券现金流映射：将付息债券分解为零息债券，使用相应的零息利率进行折现计算现值
}

\begin{question}
    The VaR percentages at the 95\% confidence level for a bond with maturities ranging from one year to five years are as follows:
    \begin{table}[H]
        \centering
        \begin{tabular}{|c|c|}
            \hline
            Maturity & VaR \% \\
            \hline
            1 & 0.5234 \\
            2 & 1.0987 \\
            3 & 1.6542 \\
            4 & 2.1985 \\
            5 & 2.7043 \\
            \hline
        \end{tabular}
    \end{table}

    A bond portfolio consists of a \$200 million bond maturing in two years and a \$200 million bond maturing in four years. What is the VaR of this bond portfolio using the principal VaR mapping method?
\end{question}

\begin{choices}
    \item \$3.308 million
    \item \$4.396 million
    \item \$6.154 million
    \item \$6.596 million
\end{choices}

\begin{correctanswer}
    D
\end{correctanswer}

\begin{explanation}
    \textbf{A选项错误。}
    \begin{itemize}
        \item 此选项错误，因为：
        \begin{itemize}
            \item 计算的VaR太低
            \item 计算错误可能出现在使用了错误的VaR百分比
        \end{itemize}
    \end{itemize}
    \textbf{B选项错误。}
    \begin{itemize}
        \item 此选项错误，因为：
        \begin{itemize}
            \item 计算的VaR太低
            \item 计算错误可能出现在使用了错误的VaR百分比
        \end{itemize}
    \end{itemize}
    \textbf{C选项错误。}
    \begin{itemize}
        \item 此选项错误，因为：
        \begin{itemize}
            \item 计算的VaR太低
            \item 计算错误可能出现在使用了错误的VaR百分比
        \end{itemize}
    \end{itemize}
    \textbf{D选项正确。}
    \begin{itemize}
        \item 此选项正确，因为：
        \begin{itemize}
            \item 步骤1：计算平均期限
            \begin{itemize}
                \item (2 + 4) / 2 = 3年
            \end{itemize}
            \item 步骤2：查找3年期债券的VaR百分比
            \begin{itemize}
                \item 从表格中：1.6542\%
            \end{itemize}
            \item 步骤3：计算投资组合VaR
            \begin{itemize}
                \item 投资组合总价值 = \$400 million
                \item VaR = \$400 million × 1.6542\% = \$6.596 million
            \end{itemize}
        \end{itemize}
    \end{itemize}
\end{explanation}
\checklist{
    掌握本金映射法计算债券组合VaR：计算平均期限，查找对应VaR百分比，乘以组合总价值
}

\begin{question}
    Which of the following statements is/are correct regarding the pricing of fixed-income securities? 
    
    I. The call price serves as the ceiling value for a callable bond when interest rates fall.
    
    II. The more frequent the time steps used for projected fixed-income security values, the more accurate the pricing.
    
    III. The Black-Scholes-Merton model is preferred for the pricing of fixed income securities because of its continuous time assumption.
\end{question}

\begin{choices}
    \item II only
    \item I and II only
    \item I and III only
    \item I, II, and III
\end{choices}

\begin{correctanswer}
    B
\end{correctanswer}

\begin{explanation}
    \textbf{A选项错误。}
    \begin{itemize}
        \item 此选项错误，因为：
        \begin{itemize}
            \item 陈述I也正确
            \item 可赎回债券在利率下降时有上限价值
        \end{itemize}
    \end{itemize}
    \textbf{B选项正确。}
    \begin{itemize}
        \item 此选项正确，因为：
        \begin{itemize}
            \item 陈述I正确：
            \begin{itemize}
                \item 可赎回债券允许发行人以设定价格赎回
                \item 这在利率下降时创造了上限价值
                \item 保护发行人免受债券价格上涨的影响
            \end{itemize}
            \item 陈述II正确：
            \begin{itemize}
                \item 更频繁的时间步长提高准确性
                \item 更好地捕捉价格变化
                \item 减少离散化误差
            \end{itemize}
            \item 陈述III错误：
            \begin{itemize}
                \item BSM模型不适用于固定收益证券
                \item 假设没有价格上限
                \item 假设利率恒定
                \item 假设波动率恒定
            \end{itemize}
        \end{itemize}
    \end{itemize}
    \textbf{C选项错误。}
    \begin{itemize}
        \item 此选项错误，因为：
        \begin{itemize}
            \item 陈述III错误
            \item BSM模型不适用于固定收益证券
        \end{itemize}
    \end{itemize}
    \textbf{D选项错误。}
    \begin{itemize}
        \item 此选项错误，因为：
        \begin{itemize}
            \item 陈述III错误
            \item BSM模型不适用于固定收益证券
        \end{itemize}
    \end{itemize}
\end{explanation}
\checklist{
    掌握固定收益证券定价：可赎回债券有上限价值，时间步长越短定价越准确，BSM模型不适用于固定收益证券
}


\begin{question}
    A fund manager owns a portfolio of options on a non-dividend paying stock DEF. The portfolio is made up of 10,000 deep in-the-money call options on DEF and 20,000 deep out-of-the-money call options on DEF. The portfolio also contains 18,000 forward contracts on DEF. Currently, DEF is trading at USD 85. Assuming 250 trading days in a year and the volatility of DEF is 20\% per year, which of the following amounts would be closest to the 1-day VaR of the portfolio at the 99\% confidence level?
\end{question}

\begin{choices}
    \item USD 35,000
    \item USD 38,500
    \item USD 41,600
    \item USD 44,800
\end{choices}

\begin{correctanswer}
    C
\end{correctanswer}

\begin{explanation}
    \textbf{A选项错误。}
    \begin{itemize}
        \item 此选项错误，因为：
        \begin{itemize}
            \item 计算的VaR太低
            \item 计算错误可能出现在delta近似上
        \end{itemize}
    \end{itemize}
    \textbf{B选项错误。}
    \begin{itemize}
        \item 此选项错误，因为：
        \begin{itemize}
            \item 计算的VaR太低
            \item 计算错误可能出现在delta近似上
        \end{itemize}
    \end{itemize}
    \textbf{C选项正确。}
    \begin{itemize}
        \item 此选项正确，因为：
        \begin{itemize}
            \item 步骤1：计算投资组合delta
            \begin{itemize}
                \item 深度实值看涨期权：delta ≈ 1
                \item 深度虚值看涨期权：delta ≈ 0
                \item 远期合约：delta = 1
                \item 总delta = 10,000 × 1 + 20,000 × 0 + 18,000 × 1 = 28,000
            \end{itemize}
            \item 步骤2：计算1天VaR
            \begin{itemize}
                \item VaR = α × S × Δ × σ × √(1/T)
                \item = 2.326 × 85 × 28,000 × 0.20 × √(1/250)
                \item = 41,600
            \end{itemize}
        \end{itemize}
    \end{itemize}
    \textbf{D选项错误。}
    \begin{itemize}
        \item 此选项错误，因为：
        \begin{itemize}
            \item 计算的VaR太高
            \item 计算错误可能出现在delta近似上
        \end{itemize}
    \end{itemize}
\end{explanation}
\checklist{
    掌握期权组合VaR的计算：使用delta-normal方法，考虑深度实值期权delta≈1，深度虚值期权delta≈0，远期delta=1
}

\begin{question}
    A hedge fund manager has to choose a risk model for a large "equity market neutral" portfolio, which has zero beta. Many of the stocks held are recent IPOs. Among the following alternatives, the best is:
\end{question}

\begin{choices}
    \item A single index model with no specific risk, estimated over the last year
    \item A diagonal index model with idiosyncratic risk, estimated over the last year
    \item A model that maps positions on industry and style factors
    \item A full covariance matrix model using a very short window
\end{choices}

\begin{correctanswer}
    C
\end{correctanswer}

\begin{explanation}
    \textbf{A选项错误。}
    \begin{itemize}
        \item 此选项错误，因为：
        \begin{itemize}
            \item 单一指数模型只考虑市场beta
            \item 投资组合构造上具有零beta
            \item 会错误地显示无风险
            \item 遗漏其他重要风险因子
        \end{itemize}
    \end{itemize}
    \textbf{B选项错误。}
    \begin{itemize}
        \item 此选项错误，因为：
        \begin{itemize}
            \item 对角模型仍然依赖历史数据
            \item IPO股票交易历史有限
            \item 会遗漏IPO特定风险
            \item 历史数据可能不可靠
        \end{itemize}
    \end{itemize}
    \textbf{C选项正确。}
    \begin{itemize}
        \item 此选项正确，因为：
        \begin{itemize}
            \item 将IPO头寸映射到行业因子
            \item 考虑风格因子（成长、价值等）
            \item 不依赖有限的历史数据
            \item 更好地捕捉IPO特定风险
            \item 更全面的风险评估
        \end{itemize}
    \end{itemize}
    \textbf{D选项错误。}
    \begin{itemize}
        \item 此选项错误，因为：
        \begin{itemize}
            \item 很短的窗口导致不可靠的估计
            \item 高估计误差
            \item 遗漏长期风险模式
            \item 不适合IPO股票
        \end{itemize}
    \end{itemize}
\end{explanation}
\checklist{
    掌握风险模型选择：对于IPO股票组合，应使用行业和风格因子映射模型，避免依赖有限的历史数据
}


\begin{question}
    A trader has an option position in natural gas with a delta of 150,000 units and gamma of -75,000 units per dollar move in price. Using the delta-gamma methodology, compute the VaR on this position, assuming the extreme move on natural gas is \$3.00 per unit.
\end{question}

\begin{choices}
    \item \$225,000
    \item \$450,000
    \item \$675,000
    \item \$900,000
\end{choices}

\begin{correctanswer}
    C
\end{correctanswer}

\begin{explanation}
    \textbf{A选项错误。}
    \begin{itemize}
        \item 此选项错误，因为：
        \begin{itemize}
            \item 计算的VaR太低
            \item 计算错误可能出现在delta-gamma近似上
        \end{itemize}
    \end{itemize}
    \textbf{B选项错误。}
    \begin{itemize}
        \item 此选项错误，因为：
        \begin{itemize}
            \item 计算的VaR太低
            \item 计算错误可能出现在delta-gamma近似上
        \end{itemize}
    \end{itemize}
    \textbf{C选项正确。}
    \begin{itemize}
        \item 此选项正确，因为：
        \begin{itemize}
            \item 使用delta-gamma近似：
            \begin{itemize}
                \item VaR(df) = Δ × VaR(dS) + (1/2)Γ × VaR(dS)²
                \item = 150,000 × (-3.00) + (1/2) × (-75,000) × (-3.00)²
                \item = -450,000 + (-337,500)
                \item = -\$675,000
            \end{itemize}
            \item 负号表示损失
        \end{itemize}
    \end{itemize}
    \textbf{D选项错误。}
    \begin{itemize}
        \item 此选项错误，因为：
        \begin{itemize}
            \item 计算的VaR太高
            \item 计算错误可能出现在delta-gamma近似上
        \end{itemize}
    \end{itemize}
\end{explanation}
\checklist{
    掌握delta-gamma VaR计算：VaR(df) = Δ × VaR(dS) + (1/2)Γ × VaR(dS)²
}

    

\begin{question}
    A trader has an option position in crude oil with a delta of 120,000 barrels and gamma of -60,000 barrels per dollar move in price. Using the delta-gamma methodology, compute the VaR on this position, assuming the extreme move on crude oil is \$2.50 per barrel.
\end{question}

\begin{choices}
    \item \$150,000
    \item \$300,000
    \item \$450,000
    \item \$600,000
\end{choices}

\begin{correctanswer}
    C
\end{correctanswer}

\begin{explanation}
    \textbf{A选项错误。}
    \begin{itemize}
        \item This option is incorrect because:
        \begin{itemize}
            \item The calculated VaR is too low
            \item The calculation error likely occurred in the delta-gamma approximation
        \end{itemize}
    \end{itemize}
    \textbf{B选项错误。}
    \begin{itemize}
        \item This option is incorrect because:
        \begin{itemize}
            \item The calculated VaR is too low
            \item The calculation error likely occurred in the delta-gamma approximation
        \end{itemize}
    \end{itemize}
    \textbf{C选项正确。}
    \begin{itemize}
        \item This option is correct because:
        \begin{itemize}
            \item Using delta-gamma approximation:
            \begin{itemize}
                \item VaR(df) = Δ × VaR(dS) + (1/2)Γ × VaR(dS)²
                \item = 120,000 × (-2.50) + (1/2) × (-60,000) × (-2.50)²
                \item = -300,000 + (-187,500)
                \item = -\$450,000
            \end{itemize}
            \item The negative sign indicates a loss
        \end{itemize}
    \end{itemize}
    \textbf{D选项错误。}
    \begin{itemize}
        \item This option is incorrect because:
        \begin{itemize}
            \item The calculated VaR is too high
            \item The calculation error likely occurred in the delta-gamma approximation
        \end{itemize}
    \end{itemize}
\end{explanation}
\checklist{
    掌握delta-gamma VaR计算
}

    
\begin{question}
    Michael Wang is a risk manager who has been assigned the task of designing a risk engine for VaR mapping. Which of the following statements accurately describes VaR mapping?
\end{question}

\begin{choices}
    \item Duration is an important factor in mapping equity portfolios.
    \item Delta-normal method is an appropriate method for estimating VaR for mapping options and interest-rate swaps.
    \item VaR mapping involves identifying common risk factors among positions in a portfolio and mapping all positions to a single risk factor.
    \item A return-based analysis may fail to spot style drift or hidden risks.
\end{choices}

\begin{correctanswer}
    D
\end{correctanswer}

\begin{explanation}
    \textbf{A选项错误。}
    \begin{itemize}
        \item 此选项错误，因为：
        \begin{itemize}
            \item 久期适用于固定收益投资组合
            \item 股票投资组合使用不同的风险因子
            \item 贝塔系数对股票更相关
        \end{itemize}
    \end{itemize}
    \textbf{B选项错误。}
    \begin{itemize}
        \item 此选项错误，因为：
        \begin{itemize}
            \item Delta-normal方法不适用于期权
            \item 期权需要完全重新定价
            \item Delta-gamma方法对期权更好
        \end{itemize}
    \end{itemize}
    \textbf{C选项错误。}
    \begin{itemize}
        \item 此选项错误，因为：
        \begin{itemize}
            \item VaR映射使用多个风险因子
            \item 不同头寸映射到不同因子
            \item 单一因子过于简化
        \end{itemize}
    \end{itemize}
    \textbf{D选项正确。}
    \begin{itemize}
        \item 此选项正确，因为：
        \begin{itemize}
            \item 基于收益的分析有局限性：
            \begin{itemize}
                \item 可能忽略风格漂移
                \item 可能忽略隐藏风险
                \item 历史收益可能不反映未来风险
            \end{itemize}
            \item VaR映射提供更好的风险评估：
            \begin{itemize}
                \item 识别共同风险因子
                \item 将头寸映射到适当的因子
                \item 更全面的风险分析
            \end{itemize}
        \end{itemize}
    \end{itemize}
\end{explanation}

\checklist{
    掌握VaR映射的特点：识别共同风险因子，将头寸映射到多个风险因子，基于收益的分析可能忽略风格漂移和隐藏风险
}

\begin{question}
    An analyst is using the delta-normal method to determine the VaR of a fixed income portfolio. The portfolio contains a long position in 1-year bonds with a \$2.5 million face value and a 5.5\% coupon that is paid semi-annually. The interest rates on six-month and twelve-month maturity zero-coupon bonds are, respectively, 3.2\% and 3.8\%. Mapping the long position to standard positions in the six-month and twelve-month zeros, respectively, provides which of the following mapped positions?
\end{question}

\begin{choices}
    \item \$68,750 and \$2,568,750
    \item \$66,234 and \$2,478,456
    \item \$67,123 and \$2,485,678
    \item \$69,456 and \$2,580,234
\end{choices}

\begin{correctanswer}
    C
\end{correctanswer}

\begin{explanation}
    \textbf{A选项错误。}
    \begin{itemize}
        \item 此选项错误，因为：
        \begin{itemize}
            \item 显示现金流但未折现
            \item 计算错误出现在未应用现值
        \end{itemize}
    \end{itemize}
    \textbf{B选项错误。}
    \begin{itemize}
        \item 此选项错误，因为：
        \begin{itemize}
            \item 计算的现值不正确
            \item 计算错误可能出现在折现过程中
        \end{itemize}
    \end{itemize}
    \textbf{C选项正确。}
    \begin{itemize}
        \item 此选项正确，因为：
        \begin{itemize}
            \item 步骤1：计算现金流
            \begin{itemize}
                \item 6个月息票 = \$2,500,000 × (0.055/2) = \$68,750
                \item 12个月支付 = \$68,750 + \$2,500,000 = \$2,568,750
            \end{itemize}
            \item 步骤2：计算现值
            \begin{itemize}
                \item 6个月零息：\$68,750 / (1 + 0.032/2) = \$67,123
                \item 12个月零息：\$2,568,750 / (1 + 0.038) = \$2,485,678
            \end{itemize}
        \end{itemize}
    \end{itemize}
    \textbf{D选项错误。}
    \begin{itemize}
        \item 此选项错误，因为：
        \begin{itemize}
            \item 计算的现值不正确
            \item 计算错误可能出现在折现过程中
        \end{itemize}
    \end{itemize}
\end{explanation}

\checklist{
    掌握债券现金流映射：将付息债券分解为零息债券，使用相应的零息利率进行折现计算现值
}

\begin{question}
    A risk analyst at an investment bank is evaluating methods for estimating portfolio VaR. The analyst considers the implications of applying a VaR mapping approach to the bank's portfolio of several fixed-income instruments and derivative contracts. Which of the following statements would the analyst be correct to make?
\end{question}

\begin{choices}
    \item When mapping for derivative contracts, exposures should only be estimated using regression analysis and cannot be directly allocated to risk factors.
    \item The choice of mapping approach involves a trade-off between preserving the market value of the positions and preserving the market risk of the positions.
    \item As the proportion of total portfolio risk explained by general risk factors increases, the proportion of total portfolio risk attributed to specific risk decreases.
    \item The portion of portfolio risk that is explained by the variance of the risk factors is the specific risk resulting from issuer-specific price movements.
\end{choices}

\begin{correctanswer}
    C
\end{correctanswer}

\begin{explanation}
    \textbf{A选项错误。}
    \begin{itemize}
        \item 此选项错误，因为：
        \begin{itemize}
            \item 衍生品可以直接映射到风险因子
            \item 回归分析不是唯一方法
            \item Delta-normal方法常用
        \end{itemize}
    \end{itemize}
    \textbf{B选项错误。}
    \begin{itemize}
        \item 此选项错误，因为：
        \begin{itemize}
            \item 映射应同时保持市场价值和风险
            \item 不需要权衡
            \item 两者可以同时实现
        \end{itemize}
    \end{itemize}
    \textbf{C选项正确。}
    \begin{itemize}
        \item 此选项正确，因为：
        \begin{itemize}
            \item 一般风险因子和特定风险是互补的
            \item 一个增加时，另一个减少
            \item 总风险保持不变
            \item 这是风险分解的基本原则
        \end{itemize}
    \end{itemize}
    \textbf{D选项错误。}
    \begin{itemize}
        \item 此选项错误，因为：
        \begin{itemize}
            \item 风险因子方差解释系统性风险
            \item 特定风险是剩余风险
            \item 陈述颠倒了关系
        \end{itemize}
    \end{itemize}
\end{explanation}

\checklist{
    掌握VaR映射原则：一般风险因子和特定风险是互补的，映射应同时保持市场价值和风险特征
}

\newpage
\end{document}